%% Journal of Intelligent Manufacturing -- Springer Nature template (sn-jnl), author-year references.
%DIF LATEXDIFF DIFFERENCE FILE
%DIF DEL old.tex    Sat Oct 24 15:36:21 2026
%DIF ADD main.tex   Sat Oct 24 15:36:21 2026
%% Line numbers for review (option lineno); remove the option for the final version.
\documentclass[pdflatex,sn-basic,lineno]{sn-jnl}

\usepackage{graphicx}
\usepackage{multirow}
\usepackage{amsmath,amssymb,amsfonts}
\usepackage{amsthm}
\usepackage{booktabs}
\usepackage{array}
\usepackage{tikz}
\usetikzlibrary{arrows.meta,positioning,calc}
\usepackage{siunitx}
\usepackage{manyfoot}   % the class's begin-document footnote hook expects this package to be loaded
\usepackage{placeins}   % \FloatBarrier keeps floats from drifting past sections
\graphicspath{{figures/}}

\raggedbottom

%% Float placement: allow figures and tables to share pages with text and with each other.
\renewcommand{\topfraction}{0.9}
\renewcommand{\bottomfraction}{0.7}
\renewcommand{\textfraction}{0.07}
\renewcommand{\floatpagefraction}{0.8}
\setcounter{topnumber}{3}
\setcounter{totalnumber}{4}
\setlength{\textfloatsep}{10pt plus 3pt minus 2pt}
\setlength{\floatsep}{10pt plus 3pt minus 2pt}
\setlength{\intextsep}{10pt plus 3pt minus 2pt}
\setlength{\belowcaptionskip}{3pt}
%% Float pages fill from the top instead of centring their floats.
\makeatletter
\setlength{\@fptop}{0pt}
\setlength{\@fpsep}{12pt plus 2fil}
\setlength{\@fpbot}{0pt plus 1fil}
\makeatother
%% Table body size: the class's own 8 bp table font (its \footnotesize is 7 pt and its \scriptsize 9 pt).
\newcommand{\tabsize}{\fontsize{8}{9.5}\selectfont}
%% Reference list: the class sets it at normal size with entries 1em apart; the 8 bp table font with
%% hanging indents keeps the entries distinct on fewer pages.
\setlength{\bibsep}{0pt}
\makeatletter
\def\bibfont{\reset@font\fontfamily{\rmdefault}\fontsize{8}{9.5}\selectfont}
\makeatother
%% DOIs in the reference list as full links (https://doi.org/...), as the journal asks.
\newcommand{\doi}[1]{\href{https://doi.org/#1}{\nolinkurl{https://doi.org/#1}}}
%% Ragged-right paragraph column for text tables.
\newcolumntype{L}[1]{>{\raggedright\arraybackslash\hspace{0pt}}p{#1}}

%% Set \anontrue to produce an anonymised copy (no author block, no identifying declarations).
\newcommand{\pending}[1]{\textit{[#1 pending]}}
%DIF 52a52-53
%% Tables of the electronic supplementary material. %DIF > 
\newcommand{\ORt}[1]{Online Resource~1, Table~S#1} %DIF > 
%DIF -------
\newif\ifanon
\anonfalse
%DIF PREAMBLE EXTENSION ADDED BY LATEXDIFF
%DIF UNDERLINE PREAMBLE %DIF PREAMBLE
\RequirePackage[normalem]{ulem} %DIF PREAMBLE
\RequirePackage{color}\definecolor{RED}{rgb}{1,0,0}\definecolor{BLUE}{rgb}{0,0,1} %DIF PREAMBLE
\providecommand{\DIFadd}[1]{{\protect\color{blue}\uwave{#1}}} %DIF PREAMBLE
\providecommand{\DIFdel}[1]{{\protect\color{red}\sout{#1}}}                      %DIF PREAMBLE
%DIF SAFE PREAMBLE %DIF PREAMBLE
\providecommand{\DIFaddbegin}{} %DIF PREAMBLE
\providecommand{\DIFaddend}{} %DIF PREAMBLE
\providecommand{\DIFdelbegin}{} %DIF PREAMBLE
\providecommand{\DIFdelend}{} %DIF PREAMBLE
\providecommand{\DIFmodbegin}{} %DIF PREAMBLE
\providecommand{\DIFmodend}{} %DIF PREAMBLE
%DIF FLOATSAFE PREAMBLE %DIF PREAMBLE
\providecommand{\DIFaddFL}[1]{\DIFadd{#1}} %DIF PREAMBLE
\providecommand{\DIFdelFL}[1]{\DIFdel{#1}} %DIF PREAMBLE
\providecommand{\DIFaddbeginFL}{} %DIF PREAMBLE
\providecommand{\DIFaddendFL}{} %DIF PREAMBLE
\providecommand{\DIFdelbeginFL}{} %DIF PREAMBLE
\providecommand{\DIFdelendFL}{} %DIF PREAMBLE
\newcommand{\DIFscaledelfig}{0.5}
%DIF HIGHLIGHTGRAPHICS PREAMBLE %DIF PREAMBLE
\RequirePackage{settobox} %DIF PREAMBLE
\RequirePackage{letltxmacro} %DIF PREAMBLE
\newsavebox{\DIFdelgraphicsbox} %DIF PREAMBLE
\newlength{\DIFdelgraphicswidth} %DIF PREAMBLE
\newlength{\DIFdelgraphicsheight} %DIF PREAMBLE
% store original definition of \includegraphics %DIF PREAMBLE
\LetLtxMacro{\DIFOincludegraphics}{\includegraphics} %DIF PREAMBLE
\newcommand{\DIFaddincludegraphics}[2][]{{\color{blue}\fbox{\DIFOincludegraphics[#1]{#2}}}} %DIF PREAMBLE
\newcommand{\DIFdelincludegraphics}[2][]{% %DIF PREAMBLE
\sbox{\DIFdelgraphicsbox}{\DIFOincludegraphics[#1]{#2}}% %DIF PREAMBLE
\settoboxwidth{\DIFdelgraphicswidth}{\DIFdelgraphicsbox} %DIF PREAMBLE
\settoboxtotalheight{\DIFdelgraphicsheight}{\DIFdelgraphicsbox} %DIF PREAMBLE
\scalebox{\DIFscaledelfig}{% %DIF PREAMBLE
\parbox[b]{\DIFdelgraphicswidth}{\usebox{\DIFdelgraphicsbox}\\[-\baselineskip] \rule{\DIFdelgraphicswidth}{0em}}\llap{\resizebox{\DIFdelgraphicswidth}{\DIFdelgraphicsheight}{% %DIF PREAMBLE
\setlength{\unitlength}{\DIFdelgraphicswidth}% %DIF PREAMBLE
\begin{picture}(1,1)% %DIF PREAMBLE
\thicklines\linethickness{2pt} %DIF PREAMBLE
{\color[rgb]{1,0,0}\put(0,0){\framebox(1,1){}}}% %DIF PREAMBLE
{\color[rgb]{1,0,0}\put(0,0){\line( 1,1){1}}}% %DIF PREAMBLE
{\color[rgb]{1,0,0}\put(0,1){\line(1,-1){1}}}% %DIF PREAMBLE
\end{picture}% %DIF PREAMBLE
}\hspace*{3pt}}} %DIF PREAMBLE
} %DIF PREAMBLE
\LetLtxMacro{\DIFOaddbegin}{\DIFaddbegin} %DIF PREAMBLE
\LetLtxMacro{\DIFOaddend}{\DIFaddend} %DIF PREAMBLE
\LetLtxMacro{\DIFOdelbegin}{\DIFdelbegin} %DIF PREAMBLE
\LetLtxMacro{\DIFOdelend}{\DIFdelend} %DIF PREAMBLE
\DeclareRobustCommand{\DIFaddbegin}{\DIFOaddbegin \let\includegraphics\DIFaddincludegraphics} %DIF PREAMBLE
\DeclareRobustCommand{\DIFaddend}{\DIFOaddend \let\includegraphics\DIFOincludegraphics} %DIF PREAMBLE
\DeclareRobustCommand{\DIFdelbegin}{\DIFOdelbegin \let\includegraphics\DIFdelincludegraphics} %DIF PREAMBLE
\DeclareRobustCommand{\DIFdelend}{\DIFOaddend \let\includegraphics\DIFOincludegraphics} %DIF PREAMBLE
\LetLtxMacro{\DIFOaddbeginFL}{\DIFaddbeginFL} %DIF PREAMBLE
\LetLtxMacro{\DIFOaddendFL}{\DIFaddendFL} %DIF PREAMBLE
\LetLtxMacro{\DIFOdelbeginFL}{\DIFdelbeginFL} %DIF PREAMBLE
\LetLtxMacro{\DIFOdelendFL}{\DIFdelendFL} %DIF PREAMBLE
\DeclareRobustCommand{\DIFaddbeginFL}{\DIFOaddbeginFL \let\includegraphics\DIFaddincludegraphics} %DIF PREAMBLE
\DeclareRobustCommand{\DIFaddendFL}{\DIFOaddendFL \let\includegraphics\DIFOincludegraphics} %DIF PREAMBLE
\DeclareRobustCommand{\DIFdelbeginFL}{\DIFOdelbeginFL \let\includegraphics\DIFdelincludegraphics} %DIF PREAMBLE
\DeclareRobustCommand{\DIFdelendFL}{\DIFOaddendFL \let\includegraphics\DIFOincludegraphics} %DIF PREAMBLE
%DIF COLORLISTINGS PREAMBLE %DIF PREAMBLE
\RequirePackage{listings} %DIF PREAMBLE
\RequirePackage{color} %DIF PREAMBLE
\lstdefinelanguage{DIFcode}{ %DIF PREAMBLE
%DIF DIFCODE_UNDERLINE %DIF PREAMBLE
  moredelim=[il][\color{red}\sout]{\%DIF\ <\ }, %DIF PREAMBLE
  moredelim=[il][\color{blue}\uwave]{\%DIF\ >\ } %DIF PREAMBLE
} %DIF PREAMBLE
\lstdefinestyle{DIFverbatimstyle}{ %DIF PREAMBLE
	language=DIFcode, %DIF PREAMBLE
	basicstyle=\ttfamily, %DIF PREAMBLE
	columns=fullflexible, %DIF PREAMBLE
	keepspaces=true %DIF PREAMBLE
} %DIF PREAMBLE
\lstnewenvironment{DIFverbatim}{\lstset{style=DIFverbatimstyle}}{} %DIF PREAMBLE
\lstnewenvironment{DIFverbatim*}{\lstset{style=DIFverbatimstyle,showspaces=true}}{} %DIF PREAMBLE
%DIF END PREAMBLE EXTENSION ADDED BY LATEXDIFF

\begin{document}

\title[When can a prediction replace a trial?]{When can a prediction replace a physical trial? A production-referenced evaluation of design-stage decision rules in deep drawing}

\ifanon
\author*[1]{\fnm{Anonymised} \sur{for review}}
\affil*[1]{\orgname{Affiliation withheld for review}}
\else
\author*[1]{\fnm{H\"useyin Tayyer} \sur{Canseven}}\email{huseyin.canseven@lut.fi}

\affil*[1]{\orgdiv{Department of Electrical Engineering, LUT School of Energy Systems}, \orgname{LUT University}, \orgaddress{\street{Yliopistonkatu 34}, \city{Lappeenranta}, \postcode{53850}, \country{Finland}}}
\fi

%% BEGIN ABSTRACT
\DIFdelbegin %DIFDELCMD < \abstract{Simulation and machine learning inform design-stage manufacturability decisions, yet they are judged by prediction error, not by whether their verdicts could replace a physical trial. This paper proposes a production-referenced evaluation of design-stage decision rules. Its unit, the within-run floor, is the difference that drift and scatter of production and measurement produce between two batches of one design within a run. For each evidence relation of a new design to produced ones (a sibling, a new process setting or a new family), a rule that returns meets, fails or trial is characterised by its decisive distance, beyond which at least 95\,\% of its verdicts are correct, and its safe distance, beyond which at most 5\,\% are wrong; a guard band turns them into single verdicts. An exploratory application to one laboratory deep-drawing process, with alternatives produced once and an uncalibrated simulation dataset, scored simulation-based, learned and production-record rules. The nearest produced setting, committed to in advance, decided at 3.4 floors with a sibling and 8.5 for a new process setting. No simulation-based or learned rule was more decisive; a Gaussian process specified for small data reduced to a force trend. Intervals calibrated on the same relation in the other family, not learning, made the production record safe for new settings. For a new family only simulation-based rules decided, none closer than 30 floors. Within a family, a production record indexed by its relation to the coming design is the baseline a learned manufacturability check has to beat.}
%DIFDELCMD < %%%
\DIFdelend \DIFaddbegin \abstract{Simulation and machine learning inform design-stage manufacturability decisions, yet they are judged by prediction error, not by whether their verdicts could replace a physical trial. This paper proposes a production-referenced evaluation of design-stage decision rules. Its unit, the within-run floor, is the difference that drift and scatter of production and measurement produce between two batches of one design within a run. For each evidence relation of a new design to produced ones (a sibling, a new process setting or a new family), a rule that returns meets, fails or trial is characterised by its decisive distance, beyond which at least 95\,\% of its verdicts are correct, and its safe distance, beyond which at most 5\,\% are wrong; a guard band turns them into single verdicts. An exploratory application to one laboratory deep-drawing process, with alternatives produced once and an uncalibrated simulation dataset, scored simulation-based, learned and production-record rules. The nearest produced setting, one rule for every relation, decided at 3.4 floors with a sibling and 8.5 for a new process setting. No simulation-based or learned rule was more decisive; a Gaussian process specified for small data reduced to a force trend. Borrowing the other family's errors on the same relation, not learning, made the production record safe for new settings. For a new family only simulation-based rules decided, none closer than 30 floors. Within a family, a production record indexed by its relation to the coming design is the baseline a learned manufacturability check has to beat.}
\DIFaddend %% END ABSTRACT

\keywords{decision support, design for manufacturing, machine learning, uncertainty quantification, selective prediction, deep drawing}

\maketitle

%% ----------------------------------------------------------------------
\section{Introduction}\label{sec:intro}

Whether a design alternative can be made to its requirements is decided long before it is made. In sheet-metal forming the decision rests on a finite-element simulation \citep{banabic2010sheet}: the team compares the predicted characteristics with the requirements and releases the alternative, changes it or sends it to a physical trial. Artificial intelligence makes this loop faster: surrogates replace simulations, models learn manufacturability from earlier and synthetic designs \citep{yan2026machine, ghadai2018learning, ballegeer2026bendfm}, and predictive-quality models and virtual metrology estimate quality from process data \citep{tercan2022machine, dreyfus2022virtual}. Such models are judged by their prediction error on held-out data, which leaves open the question the team has to answer: \emph{can this prediction be relied on to decide, or is a trial still needed?} Mechanical design cannot be verified by simulation the way integrated circuits are \citep{whitney1996why}, and the answer depends on how much production has been seen and how the new design relates to it.

Uncertainty quantification supplies intervals \citep{nemani2023uncertainty}, and virtual metrology attaches a reliance index to each prediction \citep{cheng2008evaluating}, but neither says how close to a requirement a verdict can be relied on, in a unit that production defines, for each relation of a new design to the produced evidence. This paper asks three questions. How can the fitness of a design-stage prediction to decide be expressed on a production-referenced scale, for any predictor and characteristic (RQ1)? How does it depend on the amount of production evidence and on how the new design relates to it (RQ2)? And, with six to nine produced alternatives, what do a generic simulation dataset, Gaussian processes, gradient boosting and calibrated intervals contribute against a production-record baseline (RQ3)?

The proposed framework has the \emph{within-run floor} as its unit: the difference that drift and scatter of production and measurement produce between two batches of one design within a run. A rule that turns the design-stage evidence into a verdict (\emph{meets}, \emph{fails} or a physical \emph{trial}) is scored, for each \emph{evidence relation} between a new design and the produced alternatives, by two points on its operating characteristic, its \emph{decisive} and \emph{safe distances}; a procedure turns them into a guard band for single verdicts and prices them by the costs of errors and trials. The decisions in scope are those of process planning, the approval of process changes in series production\DIFaddbegin \DIFadd{, where a customer may require a resubmission run, }\DIFaddend and the qualification of AI-based tools that support them, taken by a team that has produced variants of a design family; \DIFaddbegin \DIFadd{a new family stands for the feasibility question before }\DIFaddend the tryout of a new die, \DIFdelbegin \DIFdel{in which the tool itself}\DIFdelend \DIFaddbegin \DIFadd{the tryout itself, whose tool }\DIFaddend changes between loops, \DIFdelbegin \DIFdel{is not}\DIFdelend \DIFaddbegin \DIFadd{being out of scope}\DIFaddend .

The application uses open data: 9,000 deep-drawn and cut cups of dual-phase steel, produced as 18 design--process alternatives of 500 consecutive parts each on one laboratory tool set \citep{baum2026rddac}, with finite-element simulations of the same process from a generic dataset \citep{baum2025ddacs}. Each alternative was produced once, and the simulations were not calibrated to the process. The study is exploratory: it shows whether the framework discriminates between rules and evidence situations on real data, and its findings are hypotheses for confirmation. The contributions are the framework (Section~\ref{sec:method}); an evaluation design that separates evidence relations, isolates what the simulation adds and tests the learners against a production-record baseline (Section~\ref{sec:demo}); findings for this process (Section~\ref{sec:results}); and a proposed reporting checklist with an open scoring function (Section~\ref{sec:reuse}).

%% ----------------------------------------------------------------------
\section{Background and related work}\label{sec:background}

\subsection{Machine learning for manufacturability, predictive quality and virtual metrology}

Design for manufacturing codifies what a process can make into rules and cost models \citep{boothroyd2010product}; machine learning learns manufacturability from CAD models \citep{ghadai2018learning, zhang2018featurenet}, synthetic sheet-metal datasets \citep{ballegeer2026bendfm}, process data and simulated design spaces \citep{yan2026machine}. \DIFdelbegin \DIFdel{In production, predictive-quality }\DIFdelend \DIFaddbegin \DIFadd{Predictive-quality }\DIFaddend models and virtual metrology estimate quality from process data \citep{wuest2016machine, tercan2022machine, dreyfus2022virtual}\DIFdelbegin \DIFdel{; their uncertainty is quantified by }\DIFdelend \DIFaddbegin \DIFadd{, with uncertainty from }\DIFaddend Bayesian, ensemble and conformal methods \citep{nemani2023uncertainty, heddoub2026conformal}\DIFdelbegin \DIFdel{. Automatic virtual metrology goes further and }\DIFdelend \DIFaddbegin \DIFadd{; automatic virtual metrology }\DIFaddend attaches to every prediction a reliance index and a \DIFdelbegin \DIFdel{global }\DIFdelend similarity index that flags \DIFdelbegin \DIFdel{process data unlike the training data }\DIFdelend \DIFaddbegin \DIFadd{unfamiliar process data }\DIFaddend \citep{cheng2008evaluating, cheng2012developing}, the closest analogue \DIFdelbegin \DIFdel{in this literature }\DIFdelend to the guard band \DIFdelbegin \DIFdel{and the applicability guards }\DIFdelend of this paper. Models are carried across products and processes by transfer learning and domain adaptation \citep{maschler2021deep}, also from simulation to real parts \citep{tercan2018transfer}\DIFdelbegin \DIFdel{. In sheet forming , machine learning predicts }\DIFdelend \DIFaddbegin \DIFadd{; in sheet forming they predict }\DIFaddend defects and springback from simulated and measured data \citep{cao2024artificial, heinzelmann2025datadriven}, and discrepancies between synthetic and real \DIFaddbegin \DIFadd{deep-drawing }\DIFaddend data are being addressed \DIFdelbegin \DIFdel{for deep drawing }\DIFdelend \citep{baum2026addressing}. These models\DIFdelbegin \DIFdel{are evaluated by held-out accuracy or coverage}\DIFdelend , and language models \DIFdelbegin \DIFdel{by }\DIFdelend \DIFaddbegin \DIFadd{in }\DIFaddend case studies \citep{makatura2024how} and benchmarks \citep{picard2025concept}, \DIFaddbegin \DIFadd{are evaluated by held-out accuracy or coverage, }\DIFaddend not by their decisions against \DIFdelbegin \DIFdel{the productionof the alternatives they judge}\DIFdelend \DIFaddbegin \DIFadd{production}\DIFaddend .

\subsection{Simulation credibility and the validation of forming simulation}

Verification and validation ask whether a model is accurate enough for its intended use \DIFdelbegin \DIFdel{, distinguishing its validation from its application domain \mbox{%DIFAUXCMD
\citep{oberkampf2010verification}}\hskip0pt%DIFAUXCMD
; }\DIFdelend \DIFaddbegin \DIFadd{\mbox{%DIFAUXCMD
\citep{oberkampf2010verification}}\hskip0pt%DIFAUXCMD
, and }\DIFaddend credibility assessment scales the required evidence with the consequence of a wrong decision \citep{asme2018assessing}. Bayesian calibration estimates calibration parameters together with a model discrepancy \citep{kennedy2001bayesian, higdon2004combining}, without which \DIFdelbegin \DIFdel{calibrated parameters }\DIFdelend \DIFaddbegin \DIFadd{they }\DIFaddend are biased and over-confident \citep{brynjarsdottir2014learning, arendt2012quantification}; multi-fidelity models combine a cheap and an expensive model through a fitted scale and a discrepancy \citep{kennedy2000predicting, legratiet2014recursive}. In sheet forming, \DIFdelbegin \DIFdel{the friction law }\DIFdelend \DIFaddbegin \DIFadd{friction }\DIFaddend \citep{hol2012advanced}, the unloading and hardening behaviour \DIFdelbegin \DIFdel{that govern }\DIFdelend \DIFaddbegin \DIFadd{behind }\DIFaddend springback \citep{yoshida2002model, wagoner2013advanced} and \DIFdelbegin \DIFdel{the elastic deflection of the tools }\DIFdelend \DIFaddbegin \DIFadd{tool deflection }\DIFaddend \citep{pilthammar2018introduction} are the main sources of simulation error; robustness analysis propagates \DIFdelbegin \DIFdel{the scatter of material and process inputs \mbox{%DIFAUXCMD
\citep{wiebenga2014effect, chen2022robustness}}\hskip0pt%DIFAUXCMD
, and }\DIFdelend \DIFaddbegin \DIFadd{input scatter \mbox{%DIFAUXCMD
\citep{wiebenga2014effect, chen2022robustness}}\hskip0pt%DIFAUXCMD
, }\DIFaddend draw-in is measured and controlled in line \DIFdelbegin \DIFdel{in series production \mbox{%DIFAUXCMD
\citep{havinga2020exploiting, fischer2018implementation}}\hskip0pt%DIFAUXCMD
. Deep-drawing }\DIFdelend \DIFaddbegin \DIFadd{\mbox{%DIFAUXCMD
\citep{havinga2020exploiting, fischer2018implementation}}\hskip0pt%DIFAUXCMD
, }\DIFaddend simulation datasets support surrogate and transfer learning \citep{baum2025ddacs, heinzelmann2025comprehensive}, and the simulation-to-part deviation has been decomposed statistically \citep{baum2026statistical}\DIFdelbegin \DIFdel{; how }\DIFdelend \DIFaddbegin \DIFadd{. How }\DIFaddend close to a requirement a calibrated or learned prediction can decide remains open.

\subsection{Decisions with abstention and their evaluation}

In conformity assessment a measured value proves conformance, proves non-conformance or leaves the decision open, with guard bands set by the measurement uncertainty \citep{iso2017decision} \DIFdelbegin \DIFdel{; the }\DIFdelend \DIFaddbegin \DIFadd{and }\DIFaddend global consumer's and producer's risks \DIFaddbegin \DIFadd{that }\DIFaddend follow from a prior distribution of the measurand \citep{jcgm2012evaluation}. \DIFdelbegin \DIFdel{In }\DIFdelend \DIFaddbegin \DIFadd{The operating characteristics of }\DIFaddend acceptance sampling and selection \DIFdelbegin \DIFdel{procedures the operating characteristic gives }\DIFdelend \DIFaddbegin \DIFadd{give }\DIFaddend the probability of a decision against the true quality, with an indifference zone inside which no decision is required \citep{bechhofer1954single}. A classifier with a reject option trades errors against abstentions \citep{chow1970optimum, geifman2017selective}, as does regression with a reject option \citep{zaoui2020regression}\DIFdelbegin \DIFdel{; }\DIFdelend \DIFaddbegin \DIFadd{, and }\DIFaddend learning to defer passes difficult cases to an expert \citep{madras2018predict}. Cost-sensitive evaluation \citep{elkan2001foundations} and decision curves \citep{vickers2006decision} price the errors\DIFdelbegin \DIFdel{, and }\DIFdelend \DIFaddbegin \DIFadd{; }\DIFaddend appropriate reliance on automation requires that its reliability be visible \DIFdelbegin \DIFdel{to its users }\DIFdelend \citep{lee2004trust}, though displayed confidence can calibrate trust without improving joint accuracy \citep{zhang2020effect, zhang2021cautionary}. The decisive and safe distances are \DIFdelbegin \DIFdel{operating characteristics of this kind}\DIFdelend \DIFaddbegin \DIFadd{such operating characteristics}\DIFaddend , and the guard band is the \DIFdelbegin \DIFdel{design-stage }\DIFdelend counterpart of a conformity-assessment one, \DIFdelbegin \DIFdel{whose dependence }\DIFdelend \DIFaddbegin \DIFadd{equally dependent }\DIFaddend on a prior\DIFdelbegin \DIFdel{of requirement distances mirrors that of the global risks}\DIFdelend .

\subsection{Production variation, measurement and design decisions}

Statistical quality control compares short- and long-term variation with tolerances through capability and performance indices \DIFdelbegin \DIFdel{, with recommended minima for existing processes }\DIFdelend \citep{montgomery2019introduction}; in an initial process study an index above 1.67 meets the acceptance criteria and one between 1.33 and 1.67 may be acceptable after a review by the customer \citep{aiag2006production}, and time-dependent models describe drifting processes \citep{iso2026process}. Precision experiments define the repeatability limit, the difference that two results under repeatability conditions exceed with a probability of 5\,\% \citep{iso2025repeatability}, and intermediate precision under conditions that differ in time \citep{iso2023intermediate}; measurement systems analysis separates repeatability from bias, linearity and stability \citep{aiag2010measurement}. \DIFdelbegin \DIFdel{In engineering design, }\DIFdelend \DIFaddbegin \DIFadd{Design }\DIFaddend decisions between alternatives with uncertain outcomes \citep{hazelrigg1998framework} are taken with margins \citep{eckert2019design} or with tests whose cost is traded against that of errors \citep{thomke2001sequential, salado2018mathematical}, and what a team knows about a new alternative depends on whether it is a variant, an adaptation or an original design \citep{pahl2007engineering}, \DIFdelbegin \DIFdel{of }\DIFdelend which the evidence relations \DIFdelbegin \DIFdel{are the counterparts }\DIFdelend \DIFaddbegin \DIFadd{mirror }\DIFaddend for production evidence.

Table~\ref{tab:positioning} places the framework among these methods. It adopts the indifference zone of operating characteristics, the guard band of conformity assessment, the selective risk of abstaining predictors and grouped cross-validation \citep{roberts2017cross}. What is new is the production-defined unit in which every rule and characteristic is scored, the qualification by evidence relation, and their combination into a procedure for design-stage decisions that any predictor can enter.

\begin{table}[t]
\tabsize
\setlength{\tabcolsep}{3.5pt}
\caption{Positioning among methods that support or evaluate design-stage manufacturability decisions}\label{tab:positioning}
\begin{tabular}{@{}L{3.2cm}L{2.5cm}L{2.9cm}L{1.9cm}L{1.5cm}@{}}
\toprule
Approach & Evidence & Output & Scored against production & Abstains \\
\midrule
DFM rules and cost models \citep{boothroyd2010product} & codified process knowledge & guideline compliance, cost & no & no \\
V\&V, credibility assessment \citep{oberkampf2010verification, asme2018assessing} & physics model, validation tests & accuracy; required evidence & on test cases & can require tests \\
Bayesian calibration \citep{kennedy2001bayesian} & simulations, experiments & predictive distribution & not as decisions & implicitly \\
Conformal prediction \citep{lei2018distribution} & exchangeable calibration data & interval with coverage & coverage & yes \\
Virtual metrology with reliance index \citep{cheng2008evaluating} & process data of produced items & value, reliance and similarity index & on held-out production & flags \\
Conformity assessment, OC curves \citep{iso2017decision, bechhofer1954single} & measurement and its uncertainty & accept, reject, undecided; risks & on the measured item & yes \\
Selective classification and regression \citep{chow1970optimum, zaoui2020regression} & labelled data & value or reject; risk--coverage & on held-out data & yes \\
ML manufacturability models \citep{ghadai2018learning, yan2026machine} & historical or simulated designs & class or value & on held-out data & rarely \\
This paper & simulations, produced alternatives & verdict; distances in floors by evidence relation & yes, per verdict & yes \\
\botrule
\end{tabular}
\end{table}

%% ----------------------------------------------------------------------
\section{The framework}\label{sec:method}

\subsection{The decision problem, its scope and its loss}\label{sec:problem}

A design alternative $a$ is a combination of product geometry and process settings produced as a series of parts, over which a quality characteristic $c$ varies. A requirement is an upper limit $R$ on the characteristic, or on its absolute value for characteristics whose target is zero, that a share $p$ of the parts must meet: alternative $a$ is \emph{adequate} for $c$ when the $p$-quantile of its parts satisfies $q_p(a,c)\le R$. At the design stage $q_p$ is unknown. A decision rule $r$ maps the available evidence to an interval $[L_r, U_r]$ for $q_p$ and to a verdict
\begin{equation}
V_r = \begin{cases} \text{meets} & U_r \le R,\\ \text{fails} & L_r > R,\\ \text{trial} & \text{otherwise.} \end{cases}
\label{eq:verdict}
\end{equation}
The three verdicts mirror the decision rules of conformity assessment \citep{iso2017decision}. A rule that returns a point ($L_r=U_r$) always decides. The evidence consists of simulations of $a$ and of the production record of other alternatives, the \emph{calibration alternatives}. The verdict costs $c_\mathrm{FA}$ when a failing alternative is accepted, $c_\mathrm{FR}$ when an adequate one is rejected and $c_\mathrm{T}$ when it is sent to trial, so a rule's expected cost per verdict is
\begin{equation}
C = c_\mathrm{FA}P_\mathrm{FA} + c_\mathrm{FR}P_\mathrm{FR} + c_\mathrm{T}P_\mathrm{T},
\label{eq:cost}
\end{equation}
which the two distances of Section~\ref{sec:distances} summarise at a fixed share of errors. Three limits of scope apply: the demonstration uses upper limits on the 95th percentile ($p=0.95$, for a normal characteristic a performance index of about 0.55), lower limits following by a change of sign; verdicts are per characteristic, a design-level verdict over all characteristics being reported as an extension (Section~\ref{sec:res-robust}); and attribute outcomes such as splits and wrinkles, which need a requirement on a rate, are not covered.

\subsection{The within-run floor}\label{sec:floor}

Production evidence comes as series of consecutive parts, each cut into batches of $B$ consecutive parts summarised by a robust centre $m$ (here the 10\,\% trimmed mean). The within-run floor of characteristic $c$, the floor for short, is
\begin{equation}
F_c = Q_\alpha\bigl\{\,|m_i - m_j| : i< j \text{ batches of the same alternative}\,\bigr\},
\label{eq:floor}
\end{equation}
the $\alpha$-quantile of the difference between two batches of one design. Under a stationary normal model with a standard deviation $\sigma_\mathrm{b}$ between batch means and $\sigma_\mathrm{w}$ between parts, $F_c = 1.96\sqrt{2}\sqrt{\sigma_\mathrm{b}^2+\sigma_\mathrm{w}^2/B}$ at $\alpha=0.95$, the form of the repeatability limit \citep{iso2025repeatability, iso1994accuracy}; because it pools batch pairs up to 450 parts apart and grows with their lag, it is closer to a time-different intermediate-precision limit \citep{iso2023intermediate} of batch centres within a run. The reference protocol is $B=50$, ten batches of a 500-part run, all pairs of batches and $\alpha=0.95$; every characteristic is scored in the floor of its own geometry, and the pooled floor \DIFdelbegin \DIFdel{serves only to convert }\DIFdelend \DIFaddbegin \DIFadd{only converts }\DIFaddend distances to physical units. It rests on three assumptions: batches of one design differ only by drift and scatter within a run (A1); future runs vary within themselves as the runs it is estimated from do (A2); and the measuring system is repeatable, its resolution small against the floor (A3).

Four properties matter for its use. It has the unit of the characteristic, so distances in floors put all characteristics and predictors on a common scale, though equal distances do not mean equal nonconforming fractions. It is a property of a protocol, and $\sigma_\mathrm{b}$ and $\sigma_\mathrm{w}$ are \DIFdelbegin \DIFdel{archived }\DIFdelend \DIFaddbegin \DIFadd{reported (\ORt{4}) }\DIFaddend so that it can be recomputed. It contains whatever drifts between batches, process and measuring system alike: without a reference artefact scanned through the series it is a production-and-measurement floor, and not the shop floor. And it excludes the variation between runs, which makes it a lower bound of the variation of series production. A new family, whose floor is unknown before it is produced, is scored in the floor of the produced family. Two alternatives $a$ and $b$ are resolvable for $c$ when their effect-to-scatter ratio $\mathrm{ESR}_c(a,b) = |\bar m_a - \bar m_b|/F_c$ is at least one, $\bar m$ being the centre of all parts of an alternative; for a continuous parameter $x$ with local sensitivity $s_c=\partial\bar m/\partial x$ the minimum resolvable change is $F_c/|s_c|$.

\subsection{Decision rules}\label{sec:rules}

Eight rules carry the main results (Table~\ref{tab:rules}). The nominal simulation (M0) is the simulation at nominal sheet thickness and friction. The scaled simulation (M1s) fits $q_p = a + b\,s_0$ to the calibration alternatives by least squares, the scale term of a multi-fidelity correction without its discrepancy. The envelope rule (M2) adds the median offset of the calibration alternatives and their largest absolute residual to the top of the simulation envelope over the dataset's thicknesses and friction coefficients. The GP-corrected envelope (M5) models the offset between $q_p$ and the envelope as a Gaussian process \citep{rasmussen2006gaussian} fitted by maximum marginal likelihood, its noise variance bounded below by the variance between the calibration series of one geometry and force (the noise-variance bound). The nearest produced setting (NN) is the mean $q_p$ of the calibration alternatives at the nearest force, a capability-database look-up, and the force trend (M1sn) the least-squares line of $q_p$ in force, the counterpart of M1s without the simulation, which within a geometry is a function of the force.

The trend GP (M7n) is a learner specified for six to nine produced alternatives: a Gaussian process of $q_p$ with a linear mean in force and lubrication whose coefficients are integrated out (universal kriging), weakly informative priors on length scales and noise instead of hard bounds, and a noise term made of the known sampling variance of each $q_p$ and an estimated between-run nugget that absorbs no modelled effect. Its counterpart with the simulation (M7) is the autoregressive multi-fidelity model \citep{kennedy2000predicting}, with the simulation in the mean and a GP discrepancy\DIFaddbegin \DIFadd{; with two calibration forces its scale on the simulation acts as a force slope}\DIFaddend . Finally, intervals are separated from predictors: P+RM puts around the prediction of a rule P a relation-matched interval, the 90\,\% split-conformal quantile \citep{lei2018distribution} of P's absolute errors in floors over the cases of the same relation in the other family, a group-conditional construction \citep{vovk2005algorithmic}. The errors are borrowed because a family's own calibration set \DIFdelbegin \DIFdel{has no case of its new-setting relation, so the intervals guarantee coverage only if the two families' errors are exchangeable, which the data cannot establish \mbox{%DIFAUXCMD
\citep{barber2023conformal}}\hskip0pt%DIFAUXCMD
; }\DIFdelend \DIFaddbegin \DIFadd{lacks the relation. The 21 (interpolated), 42 (extrapolated) or 63 cases behind a quantile come from one family and from one or a few shared fits, so they are not exchangeable with the other family's cases and the interval carries no coverage guarantee \mbox{%DIFAUXCMD
\citep{barber2023conformal}}\hskip0pt%DIFAUXCMD
; its 90\,\% is nominal. No interval from $n$ exchangeable cases promises a miss rate below $1/(n+1)$: a guaranteed 5\,\% needs 19 produced alternatives in a relation. }\DIFaddend NN+RM is NN with such an interval.

\DIFdelbegin \DIFdel{Ten }\DIFdelend \DIFaddbegin \DIFadd{Besides M7, ten }\DIFaddend further rules complete the comparison. M1 adds the median offset of the calibration alternatives to M0, and M1n is their median $q_p$; M2n and M5n are M2 and M5 without the simulation. Mc fits one friction coefficient per lubrication pattern to the draw-ins by least squares, without a discrepancy, and Mcd calibrates the friction with a GP discrepancy \citep{kennedy2001bayesian} and predicts by the posterior mixture over the friction grid. M3 corrects the matched simulation of every part by gradient boosting \DIFaddbegin \DIFadd{\mbox{%DIFAUXCMD
\citep{friedman2001greedy}}\hskip0pt%DIFAUXCMD
}\DIFaddend , and M3n models the characteristic by boosting without the simulation, both with jackknife+ intervals \DIFaddbegin \DIFadd{\mbox{%DIFAUXCMD
\citep{barber2021predictive}}\hskip0pt%DIFAUXCMD
}\DIFaddend ; M0w is the range of the simulation over thickness and friction; and M6, a robustness variant, puts a conformal margin on standardised residuals around the centre of M5.

\begin{table}[t]
\tabsize
\setlength{\tabcolsep}{3.5pt}
\caption{The rules of the main results; the further rules are defined in Section~\ref{sec:rules}}\label{tab:rules}
\DIFdelbeginFL %DIFDELCMD < \begin{tabular}{@{}L{3.6cm}L{4.6cm}L{2.3cm}L{1.4cm}@{}}
%DIFDELCMD < \toprule
%DIFDELCMD < Rule & Prediction of $q_p$ & Interval (level) & Evidence \\
%DIFDELCMD < \midrule
%DIFDELCMD < M0 nominal simulation & simulation at nominal thickness and friction & point & none \\
%DIFDELCMD < M1s scaled simulation$^\dagger$ & $a+b\,$M0, least squares & point & $q_p$ \\
%DIFDELCMD < M2 envelope + residual margin & envelope top plus the median offset & largest residual ($\frac{n-1}{n+1}$) & $q_p$ \\
%DIFDELCMD < M5 GP-corrected envelope & envelope plus a GP offset, ML-II & GP (nominal 0.9) & $q_p$ \\
%DIFDELCMD < NN nearest produced setting$^\dagger$ & mean $q_p$ at the nearest produced force & point & $q_p$ \\
%DIFDELCMD < M1sn force trend$^\dagger$ & $a+b\,$force, least squares & point & $q_p$ \\
%DIFDELCMD < M7n trend GP$^\ddagger$ & linear mean plus GP, MAP under priors & GP (nominal 0.9) & $q_p$ \\
%DIFDELCMD < NN+RM$^\ddagger$ & NN & relation-matched (0.9) & $q_p$ \\
%DIFDELCMD < \botrule
%DIFDELCMD < \end{tabular}
%DIFDELCMD < %%%
\DIFdelendFL \DIFaddbeginFL \begin{tabular}{@{}L{3.6cm}L{4.6cm}L{2.3cm}L{1.4cm}@{}}
\toprule
Rule & Prediction of $q_p$ & Interval (level) & Evidence \\
\midrule
M0 nominal simulation & simulation at nominal thickness and friction & point & none \\
M1s scaled simulation$^\dagger$ & $a+b\,$M0, least squares & point & $q_p$ \\
M2 envelope + residual margin & envelope top plus the median offset & largest residual ($\frac{n-1}{n+1}$) & $q_p$ \\
M5 GP-corrected envelope & envelope plus a GP offset, ML-II & GP (nominal 0.9) & $q_p$ \\
NN nearest produced setting$^\dagger$ & mean $q_p$ at the nearest produced force & point & $q_p$ \\
M1sn force trend$^\dagger$ & $a+b\,$force, least squares & point & $q_p$ \\
M7n trend GP$^\ddagger$ & linear mean plus GP, MAP under priors & GP (nominal 0.9) & $q_p$ \\
NN+RM$^\ddagger$ & NN & relation-matched (nominal 0.9) & $q_p$ \\
\botrule
\end{tabular}
\DIFaddendFL \footnotetext{Fitted on the calibration alternatives; descriptors: geometry indicator, blank-holder force and lubrication rank. Pre-specified unless marked: $^\dagger$ added after the first results; $^\ddagger$ added \DIFaddbeginFL \DIFaddFL{in the revisions }\DIFaddendFL for this \DIFdelbeginFL \DIFdelFL{version}\DIFdelendFL \DIFaddbeginFL \DIFaddFL{journal}\DIFaddendFL . Levels: guaranteed (largest residual, $n$ calibration alternatives: 8 with a sibling, 6 for a new setting, 9 for a new family) or nominal.}
\end{table}

\subsection{Operating characteristics: decisive and safe distances}\label{sec:distances}

A rule is scored by judging each produced alternative as a new design at requirements placed at a known distance from the truth,
\begin{equation}
R = q_p(a,c) + d\,F_c ,
\label{eq:requirement}
\end{equation}
so that $a$ is adequate exactly when $d\ge 0$. Let $\kappa(\delta)$ and $\omega(\delta)$ be the shares of correct and of wrong verdicts among all feasible verdicts at $|d|=\delta$: every case gives one verdict on an adequate alternative and, where the requirement is not negative, one on a failing alternative. The \emph{decisive distance} $\delta_\mathrm{d}$ is the smallest $\delta$ beyond which $\kappa\ge 0.95$, and the \emph{safe distance} $\delta_\mathrm{s}$ the smallest $\delta$ beyond which $\omega\le 0.05$ (Fig.~\ref{fig:schematic}a). In plain terms, for a rule that always decides, $\delta_\mathrm{d}=\delta_\mathrm{s}$ is close to the 90th percentile of its absolute error in floors; an interval rule trades decisiveness for safety, and between $\delta_\mathrm{s}$ and $\delta_\mathrm{d}$ it sends designs to trial rather than misleading. Both are points on operating characteristics over a population of design situations; $\delta_\mathrm{d}$ is the analogue of the indifference zone of a selection procedure \citep{bechhofer1954single} at a probability of correct selection of 0.95.

The distances are exact up to \DIFdelbegin \DIFdel{the }\DIFdelend \DIFaddbegin \DIFadd{their }\DIFaddend 0.05-floor grid\DIFdelbegin \DIFdel{on which they are evaluated, because $\kappa$ and $\omega$ change only where $\delta$ crosses the distance by which a case's interval excludes its truth}\DIFdelend . They are order statistics, so a few verdicts far out can move them: with 126 cases, some 240 verdicts, the 95\,\% threshold rests on about a dozen. Distances below about one floor are reported as ``$<$1'', since production reproduces the decided quantile only to about one floor within a run (Section~\ref{sec:res-floor}). The distances weigh false accepts and false rejects equally, and as $\delta$ grows, failing-side requirements become infeasible for characteristics close to zero, so the failing side carries half the weight up to a few floors and a third at about 100 floors; the safe distances \DIFdelbegin \DIFdel{on the false-accept and the false-reject }\DIFdelend \DIFaddbegin \DIFadd{of each }\DIFaddend side are therefore reported \DIFdelbegin \DIFdel{as well}\DIFdelend \DIFaddbegin \DIFadd{too}\DIFaddend .

\begin{figure}[t]
\centering
\includegraphics[width=\textwidth]{F_schematic}
\caption{Operating characteristics of one rule (M7n, extrapolated setting): verdict shares against the requirement distance, with $\delta_\mathrm{d}$ and $\delta_\mathrm{s}$ (a), and correct share among decided verdicts against the margin, with $g$ (b)}\label{fig:schematic}
\end{figure}

\subsection{Evidence relation and calibration budget}\label{sec:budget}

How a new design relates to the produced alternatives is part of the decision situation. With a \emph{sibling}, an alternative at the same process setting has been produced that differs only in the factor varied at a fixed setting (here the lubrication pattern, which the nominal simulation does not represent); a \emph{new process setting} lies between produced settings (interpolation) or outside them (extrapolation); a \emph{new family} shares no produced alternative. A team knows, before it decides, whether the produced siblings of a new variant differ resolvably from each other, so the sibling relation is stratified by this decision-time information, not by the new design's own production. The calibration budget is the dependence of the distances on the number $k$ of produced alternatives.
\DIFdelbegin \DIFdel{An applicability guard flags a design outside the calibration, in the spirit of a valid input domain \mbox{%DIFAUXCMD
\citep{malak2010using}}\hskip0pt%DIFAUXCMD
: its geometry is not calibrated, a descriptor lies outside the calibrated range, or its simulated $q_p$ lies outside that of the calibration alternatives.
}\DIFdelend 

\subsection{Procedure}\label{sec:procedure}

Figure~\ref{fig:procedure} turns the framework into six steps, steps 1 to 5 once per design family and repeated when production evidence is added, and step 6 for every new design. Steps 1 to 4 frame the decision (family, alternatives, characteristics, form of the requirement, costs), measure the floor with a stated protocol, state which differences production resolves, and calibrate the candidate rules with and without the simulation.

Step 5 qualifies the rules of the family by their distances for each evidence relation and fixes for each rule and relation a \emph{guard band} $g$: the smallest observable margin between the predicted interval and the requirement beyond which at least 95\,\% of the decided verdicts are correct. \DIFdelbegin \DIFdel{Being conditional }\DIFdelend \DIFaddbegin \DIFadd{Conditional }\DIFaddend on what is observed, $g$ depends on the distribution of requirement distances \DIFdelbegin \DIFdel{that }\DIFdelend the team expects, its prior, as the global risks of conformity assessment do \citep{jcgm2012evaluation}. The reference prior weights the 47 distances of the evaluation grid within $\pm20$ floors equally\DIFdelbegin \DIFdel{(0, $\pm0.5$ to $\pm10$ in steps of 0.5, $\pm12$, $\pm15$, $\pm20$)}\DIFdelend , requirements set a few to some tens of floors from the true 95th percentile ($\pm20$ floors are about $\pm1$\,mm of mid-side draw-in, $\pm1.6^\circ$ of wall angle, $\pm0.2$\,mm of cup depth and $\pm0.07$--0.16\,mm of dome); a release-like prior, the grid within $\pm5$ floors, stands for requirements near capability limits, and a team replaces both by the distribution of its own requirements, for example the margins between past drawing limits and achieved 95th percentiles. The band is estimated per family \DIFdelbegin \DIFdel{and validated out of sample, leaving out one alternative and, for new settings, }\DIFdelend \DIFaddbegin \DIFadd{from the relations its production exercises and validated leaving }\DIFaddend one \DIFdelbegin \DIFdel{force cell}\DIFdelend \DIFaddbegin \DIFadd{alternative out; a relation the family's production cannot exercise (interpolation in a family produced at two forces) borrows the band of the same relation in another family, with the caveat of Section~\ref{sec:rules}, or sends the design to trial. Stated costs may only widen the band, to the one minimising Eq.~(\ref{eq:cost}) over the prior, so they never lower the share correct below 95\,\%}\DIFaddend .

Step 6 establishes the relation of the new design, applies the rule qualified for it \DIFaddbegin \DIFadd{(with stated costs, the cheapest)}\DIFaddend , widens the rule's interval by $g$ floors and returns meets or fails with the margin in floors, or a trial when no rule is qualified or the requirement lies inside the widened interval. In pseudocode, step 5 is: for each rule and relation, collect the held-out verdicts over the prior; for $g = 0, 0.25, \ldots, 20$, compute the share correct among verdicts decided with a margin of at least $g$; take the smallest $g$ beyond which it stays at or above 0.95 \DIFaddbegin \DIFadd{and, with stated costs, raise it to the $g$ that minimises Eq.~(\ref{eq:cost})}\DIFaddend . Step 6 is: $V = \text{meets}$ if $U+gF\le R$, fails if $L-gF>R$, trial otherwise.
\DIFdelbegin \DIFdel{The band of the relation takes the place of the range and family guards, which here coincide with the extrapolated and new-family relations.
}\DIFdelend 

\begin{figure}[t]
\centering
\input{fig_procedure}
\caption{The procedure: steps 1--5 once per design family, step 6 for every new design; dashed boxes are the evidence and inputs it uses}\label{fig:procedure}
\end{figure}

\subsection{A worked example}\label{sec:example}

A process planner has to decide whether the concave cup \DIFaddbegin \DIFadd{of Section~\ref{sec:data} }\DIFaddend can be run at 300\,kN with medium lubrication, a setting not yet produced, without a trial run, the six alternatives at 100 and 500\,kN having been produced. Step 1 frames the decision: the characteristic is the mid-side draw-in, by which the blank-holder force is set, the requirement (chosen for illustration) that 95\,\% of the parts draw in by at most 15.30\,mm, and a trial costs a short press run, about a twentieth of a false accept. Step 2 gives the floor from the \DIFaddbegin \DIFadd{six }\DIFaddend produced series, \DIFdelbegin \DIFdel{0.050}\DIFdelend \DIFaddbegin \DIFadd{0.041\,mm; the draw-in is quantised in steps of 0.019}\DIFaddend \,mm, \DIFdelbegin \DIFdel{and step }\DIFdelend \DIFaddbegin \DIFadd{three values holding a median 89\,\% of a series. Step }\DIFaddend 3 shows that \DIFdelbegin \DIFdel{force changes of 21--67}\DIFdelend \DIFaddbegin \DIFadd{between 100 and 500\,kN a force change of about 40}\DIFaddend \,kN \DIFdelbegin \DIFdel{are }\DIFdelend \DIFaddbegin \DIFadd{is }\DIFaddend resolvable. Step 4 calibrates the rules on the six alternatives. Step 5 \DIFdelbegin \DIFdel{qualifies }\DIFdelend \DIFaddbegin \DIFadd{cannot qualify a rule for interpolation on this family, whose production holds no interpolated case, so it borrows from the convex family, produced at all three forces, assuming that the two families err alike: there }\DIFaddend the nearest produced setting \DIFdelbegin \DIFdel{for interpolation: it decides at 4.7 }\DIFdelend \DIFaddbegin \DIFadd{decides at 5.3 }\DIFaddend floors and needs a \DIFdelbegin \DIFdel{guard band of 1.5 floorson this family (2.0 pooled over the families); out of sample, 95\,\% of its decided verdicts stayed correct}\DIFdelend \DIFaddbegin \DIFadd{band of 2.5 floors, which the stated trial cost widens to 5.0}\DIFaddend . Step 6: the nearest setting, the mean of the six alternatives at the neighbouring forces, predicts 14.99\,mm, a margin of \DIFdelbegin \DIFdel{6.2 }\DIFdelend \DIFaddbegin \DIFadd{7.4 }\DIFaddend floors, more than \DIFdelbegin \DIFdel{its }\DIFdelend \DIFaddbegin \DIFadd{the }\DIFaddend band, so the planner releases the setting. The trend GP agrees with a margin of \DIFdelbegin \DIFdel{4.5 floors; }\DIFdelend \DIFaddbegin \DIFadd{5.4 floors; the intervals of }\DIFaddend M2, M5 and M5n \DIFdelbegin \DIFdel{, widened by their bands, }\DIFdelend contain the limit and would send the design to trial, and the nominal simulation, qualified for no relation, predicts 15.66\,mm. The 500 parts had a 95th percentile of 14.88\,mm, \DIFdelbegin \DIFdel{8.4 }\DIFdelend \DIFaddbegin \DIFadd{10.0 }\DIFaddend floors below the limit\DIFaddbegin \DIFadd{; Section~\ref{sec:res-procedure} reports how often such verdicts are right}\DIFaddend .

%% ----------------------------------------------------------------------
\section{Demonstration on deep drawing}\label{sec:demo}

\subsection{Alternatives and production data}\label{sec:data}

The production data are the Real Deep Drawing and Cutting dataset (RDDAC) of the University of Stuttgart \citep{baum2026rddac, baum2026statistical}. Square blanks of DP600 dual-phase steel, \DIFaddbegin \DIFadd{about 1\,mm thick and taken as }\DIFaddend 210\,mm wide\DIFdelbegin \DIFdel{and about }\DIFdelend \DIFaddbegin \DIFadd{, the blank of the simulation dataset (}\DIFaddend 1\,mm \DIFdelbegin \DIFdel{thick}\DIFdelend \DIFaddbegin \DIFadd{of width moves the measured draw-in by 0.5\,mm, 10 floors)}\DIFaddend , were deep drawn into cups (OP10) and cut (OP20) on an industrial press with one tool set, whose two die geometries, concave and convex side walls, differ in base curvature, wall angle and bottom radius at once. Each geometry was run at three blank-holder forces (100, 300 and 500\,kN) and three lubrication patterns applied to the blank (coarse, medium and fine, with series medians of 1.15--1.48, 0.92--1.23 and 0.85--1.13\,g/m$^2$ of oil): 18 alternatives, each produced as one series of 500 consecutive parts. Every part has \DIFdelbegin \DIFdel{the signals of four load cells around the die }\DIFdelend \DIFaddbegin \DIFadd{four load-cell signals }\DIFaddend at 300\,Hz, the punch position and temperature, sheet-thickness and oil-film traverses and laser scans after both operations. The series warm up by 1.0--5.0\,K over their first 150 parts from punch temperatures of 20.5--29.6\,$^\circ$C, their median sheet thickness is 0.986--0.993\,mm, and the 100\,kN series ran at 47\,mm/s against 78--79\,mm/s, so the extrapolation to 100\,kN also changes the sliding speed and stroke rate on which the lubricant film depends. The experiments were collected between September and November 2025, in contiguous blocks of 500 per alternative; the documentation states neither the order of the blocks nor whether the parts were scanned in production order, and gives no \DIFaddbegin \DIFadd{blank size, }\DIFaddend material certificate, coil record, tool drawing, press or cushion type, or \DIFdelbegin \DIFdel{method of applying the lubrication patterns}\DIFdelend \DIFaddbegin \DIFadd{how the lubrication was applied}\DIFaddend .

\subsection{Simulations}\label{sec:sims}

The simulations come from the Deep Drawing and Cutting Simulations dataset (DDACS) \citep{baum2025ddacs, heinzelmann2025comprehensive}, a generic dataset of some 32,000 LS-DYNA simulations of the same cup family that was neither set up for nor validated against the produced process. Its documentation gives rigid tools\DIFdelbegin \DIFdel{(blank holder, die, punch)}\DIFdelend , shell elements with seven integration points through the thickness and nodes about 1\,mm apart, a scaled DP600 flow curve, constant Coulomb friction, a drawing depth of 30\,mm, trimming in OP20 without modelled cutting tools, and springback after both operations; it leaves open the element formulation, whether contact accounts for the shell thickness, the yield criterion and hardening law, the blank-holder actuation and the springback procedure, and it does not represent the stroke speed. For the RDDAC tools, 396 simulations vary the sheet thickness (0.95--1.00\,mm) and the friction coefficient (0.05--0.15) at each geometry and force, all with the nominal flow curve. A part is matched to the simulation with the nearest thickness and the friction coefficient that the dataset's linear mapping assigns to its oil film, an assumption of the dataset's authors on which M3, M3n and the matched simulations depend. The nominal simulation is at 0.99\,mm, nearest the median thickness, and friction 0.10, the middle of the grid\DIFdelbegin \DIFdel{; Mc and Mcd calibrate the friction }\DIFdelend \DIFaddbegin \DIFadd{. The simulated characteristics carry extraction noise: the two directions of the quarter model differ by 1.8--2.3 floors in wall angle, and the dome is rough along the friction grid by 2.4--9.8 floors (standard deviations; \ORt{3})}\DIFaddend .

\subsection{Measurement and quality characteristics}\label{sec:qcs}

The parts were measured with a single-sided laser line scanner, which records a height and a luminescence image of 3200 by 2000 points, 0.077\,mm apart along the laser line ($x$) and 0.158\,mm apart between lines in the direction of motion ($y$), with height steps of 7.75\,\textmu m, the cup bottom facing the scanner (Fig.~\ref{fig:parts}). The scales are not traced to the part: the lateral one comes from the scanner specification\DIFdelbegin \DIFdel{and a 30}\DIFdelend \DIFaddbegin \DIFadd{, checked against a 9}\DIFaddend \,mm \DIFdelbegin \DIFdel{block, the line spacing from the scans, partly fitted to the simulation, and the height from a }\DIFdelend \DIFaddbegin \DIFadd{wide calibration block that also gives the height (}\DIFaddend 1\,mm\DIFdelbegin \DIFdel{block}\DIFdelend \DIFaddbegin \DIFadd{)}\DIFaddend , against a nominal repeatability of 5\,\textmu m\DIFaddbegin \DIFadd{; the line spacing was estimated from the scans and corrected by a fit of the simulation's scale, which the corner draw-in, the arm angle and the dome use and the mid-side draw-in, the wall angle and the depth do not}\DIFaddend . Absolute values are therefore uncalibrated, which limits every absolute comparison with the simulation \DIFaddbegin \DIFadd{(a lateral scale error of 1\,\% moves the mid-side draw-in by 0.9\,mm, 18 floors) }\DIFaddend but not the floor, a within-alternative difference that a constant scale error only rescales.

Seven characteristics were extracted with the same definitions from the scans and from the simulated nodes, referred to the plane of the cup bottom (Table~\ref{tab:qcs}): the mid-side and corner draw-in and the flange waviness, process indicators on the flange that the cut removes; the wall angle and the cup depth, measured after drawing because the cut removes the flange that defines the depth and the cut walls facing the scanner read several degrees too steep; and the arm angle and bottom dome of the cut part, the last two with requirements on absolute values. The four sides of a part measure several of them twice, which exposes direction-dependent systematic errors of up to 3.5\,mm in flange width and 10$^\circ$ in wall angle (Table~\ref{tab:qcs}); the directions kept avoid the worst of them, but their residual bias is unknown. The mid-side draw-in is quantised in steps of 0.019\,mm (0.38 floors), so \DIFdelbegin \DIFdel{that }\DIFdelend a series takes only 6 to 21 \DIFdelbegin \DIFdel{distinct }\DIFdelend values (Fig.~\ref{fig:batches}).

%% BEGIN GENERATED qcs
\begin{table}[t]
\tabsize
\setlength{\tabcolsep}{2.4pt}
\caption{The quality characteristics and their measurement by the laser scanner}\label{tab:qcs}
\begin{tabular}{@{}L{1.75cm}L{0.75cm}L{3.15cm}rrL{3.85cm}@{}}
\toprule
Characteristic & Op. & Region; definition & Step$/F$ & Values & Direction-dependent differences \\
\midrule
Draw-in, mid-side & OP10 & flange mid-sides; $(210-W_x)/2$, $W_x$ along the laser line & 0.38 & 6--21 & $W_y-W_x$ 2.29--3.47\,mm; $W_y$, in the direction of motion, scatters 3.7--17.6 times more \\
Draw-in, corner & OP10 & flange corners; mean of the two diagonals & $<$0.01 & 491--500 & the two diagonals differ by 1.64--2.17\,mm \\
Flange waviness & OP10 & flange; residual SD about a quadratic surface & $<$0.01 & 469--494 & -- \\
Wall angle & OP10 & side walls at mid-depth; mean of east and west & $<$0.01 & 497--500 & convex: north wall $-$1.64 to $-$0.92$^\circ$ from east/west (concave $-$0.05 to +0.39$^\circ$); south wall shadowed, 6.7--10.0$^\circ$ off; after the cut, north 3.3--5.6$^\circ$ off (convex) \\
Arm angle, cut & OP20 & cut arms, centre lines; mean of the four arms & $<$0.01 & 496--500 & -- \\
Cup depth & OP10 & flange above the bottom plane; median & $<$0.01 & 486--497 & the flange plane is cut away at OP20 \\
Bottom dome, cut & OP20 & bottom sag at $r=40$\,mm; paraboloid fit & $<$0.01 & 474--497 & -- \\
\botrule
\end{tabular}
\footnotetext{Op.: operation measured (OP10 after drawing, OP20 after cutting). Step$/F$: median step between adjacent values in one series over the floor; Values: distinct values per series. Differences from the redundancy of the four sides, ranges of the series medians; the directions used avoid them, their residual bias is unknown. Absolute values are not calibrated against a reference at the size of the part.}
\end{table}
%% END GENERATED qcs

\begin{figure}[t]
\centering
\includegraphics[width=\textwidth]{F_parts}
\caption{Scanned parts after drawing and cutting, depth below the cup bottom (a, b, d, e), and sections (c, f); lines and boxes: directions and regions of the characteristics}\label{fig:parts}
\end{figure}

\subsection{Evaluation design}\label{sec:evaldesign}

Every produced alternative is judged as a new design under grouped cross-validation \citep{roberts2017cross}, one scheme per evidence relation. \emph{New variant}: the alternative alone is held out and the other eight of its geometry calibrate, two siblings at its force among them. \emph{New process setting}: the three alternatives of one force of a geometry are held out, the six others calibrate; 300\,kN is interpolated, 100 and 500\,kN extrapolated. \emph{New family}: one geometry calibrates the rules for the other (two transfers); a few-shot variant adds one to three produced alternatives of the new family. Every calibrated rule is evaluated with and without the simulation (M1s/M1sn, M2/M2n, M3/M3n, M5/M5n, M7/M7n). Each relation has 126 cases (18 alternatives, 7 characteristics), but the units of replication are the 18 alternatives, whose characteristics are correlated.

The \DIFdelbegin \DIFdel{uncertainty of the distances is described in three ways. The }\DIFdelend intervals of Table~\ref{tab:distances} resample the held-out alternatives within each geometry 1000 times (those held out together for a new setting individually), each resample with a block-bootstrap replicate of every true $q_p$ (non-overlapping blocks of 50 parts), the fitted rules held fixed; common to all rules, they give paired intervals of differences, descriptive and not adjusted for the many comparisons. \DIFdelbegin \DIFdel{A refitting bootstrap redraws the three lubrication patterns of each geometry with replacement and refits every alternative-level rule, which shows the effect of losing a pattern rather than sampling uncertainty. }\DIFdelend Coverage gets cluster-bootstrap intervals over alternatives \DIFaddbegin \DIFadd{where each has its own calibration set}\DIFaddend . Each truth is the 95th percentile of a single run, so a sensitivity analysis adds to every truth a between-run component of standard deviation $\tau$ floors, from near-replicate reproducibility up to the variation between series. The selection of rules and bands is cross-fitted between the geometries, chosen on one and evaluated on the other; like any cross-validated estimate, the distances describe the fitting procedure on this design rather than one fitted rule \citep{bates2024cross}.

\subsection{Computation and development record}\label{sec:computation}

The constants of the floor and of the requirement grid were set before the decision results were seen, which the archive cannot verify \DIFdelbegin \DIFdel{, as no timestamped snapshotwas deposited at the time}\DIFdelend \DIFaddbegin \DIFadd{for want of a timestamped snapshot}\DIFaddend . The rules were not fixed in advance: M0--M5 and M3n were defined first, \DIFdelbegin \DIFdel{and everything else was added }\DIFdelend \DIFaddbegin \DIFadd{everything else }\DIFaddend after the first results \DIFaddbegin \DIFadd{(Table~\ref{tab:rules} marks the main rules)}\DIFaddend , most of it in response to reviews of earlier versions simulated with a large language model in \DIFdelbegin \DIFdel{four rounds(M0w, Mc, M1n, NN, M2n, M5n, the grouped schemes, the exact estimator, the noise-variance bound, M1s, M1sn, the jackknife+ intervals, the guard bands with their priors and the new family scored in the produced family's floor). A fifth simulated round, of the version prepared for this journal, led to M7n}\DIFdelend \DIFaddbegin \DIFadd{six rounds}\DIFaddend , \DIFdelbegin \DIFdel{M7, Mcd, the relation-matched intervals, the decision-time strata, the per-family and out-of-sample guard bands, the costs, the between-run sensitivity, the cross-fitting, the few-shot transfer and the design-level verdicts. }\DIFdelend \DIFaddbegin \DIFadd{listed in Online Resource~1. }\DIFaddend The comparison of rules is therefore exploratory and carries the optimism of a selection made on the evaluation data \citep{cawley2010over, varma2006bias}, which the cross-fitting \DIFdelbegin \DIFdel{bounds}\DIFdelend \DIFaddbegin \DIFadd{estimates}\DIFaddend ; every variant tried is reported \DIFdelbegin \DIFdel{in the paper }\DIFdelend \DIFaddbegin \DIFadd{here, in Online Resource~1 }\DIFaddend or in the \DIFdelbegin \DIFdel{archived results \mbox{%DIFAUXCMD
\citep{canseven2026archive}}\hskip0pt%DIFAUXCMD
, in which }\texttt{\DIFdel{results/findings.md}} %DIFAUXCMD
\DIFdel{maps every number of the paper to its script and result file}\DIFdelend \DIFaddbegin \DIFadd{archive \mbox{%DIFAUXCMD
\citep{canseven2026archive}}\hskip0pt%DIFAUXCMD
}\DIFaddend .

\emph{Computation.} M5 and M5n are Gaussian-process regressions \DIFdelbegin \DIFdel{(constant times }\DIFdelend \DIFaddbegin \DIFadd{with an }\DIFaddend anisotropic squared-exponential kernel \DIFdelbegin \DIFdel{plus }\DIFdelend \DIFaddbegin \DIFadd{and }\DIFaddend white noise, \DIFdelbegin \DIFdel{length scales in $[0.1, 100]$, }\DIFdelend \DIFaddbegin \DIFadd{fitted by }\DIFaddend maximum marginal likelihood \DIFdelbegin \DIFdel{) }\DIFdelend on a geometry indicator, the force in 100\,kN and the lubrication rank\DIFdelbegin \DIFdel{; the noise-variance bound is the larger of the mean sampling variance of the calibration truths and the pooled variance between their series of one geometry and force}\DIFdelend . M7n\DIFdelbegin \DIFdel{and }\DIFdelend \DIFaddbegin \DIFadd{, }\DIFaddend M7 \DIFaddbegin \DIFadd{and Mcd }\DIFaddend work in pooled floors \DIFdelbegin \DIFdel{, with }\DIFdelend \DIFaddbegin \DIFadd{under }\DIFaddend log-normal priors \DIFdelbegin \DIFdel{whose medians are }\DIFdelend \DIFaddbegin \DIFadd{with medians of }\DIFaddend 200\,kN and one pattern step for the length scales \DIFdelbegin \DIFdel{, two }\DIFdelend \DIFaddbegin \DIFadd{and two and one }\DIFaddend floors for the GP and \DIFdelbegin \DIFdel{one floor for the nugget standard deviation (log-scale standard deviations1, 1.5 and 1). Mcd has }\DIFdelend \DIFaddbegin \DIFadd{nugget standard deviations, Mcd with }\DIFaddend a uniform prior over the 11 friction values of the grid\DIFdelbegin \DIFdel{, a GP discrepancy with the same priors and a posterior from the two draw-ins}\DIFdelend . M3 and M3n use histogram gradient boosting \DIFdelbegin \DIFdel{on geometry, force, lubrication, sheet thickness and oil film, }\DIFdelend with 500 incoming conditions drawn from calibration parts of the same pattern. Software: Python 3.11, SciPy and scikit-learn \citep{pedregosa2011scikit}\DIFaddbegin \DIFadd{; every setting is in the archive}\DIFaddend .

%% ----------------------------------------------------------------------
\section{Results}\label{sec:results}

\subsection{What the floor contains}\label{sec:res-floor}

Two batches of 50 consecutive parts of one alternative differ by up to 0.051\,mm in mid-side draw-in and 0.078$^\circ$ in wall angle in 95\,\% of batch pairs (Table~\ref{tab:floors}), so one floor is 0.7 to 1.7 long-term standard deviations of the parts. The floors are set by drift within the run (Fig.~\ref{fig:batches}): without the time order they shrink by a factor of 2.1 to 4.5, and batches one position apart differ by 0.51--0.80 floors, nine apart by 1.1--1.5. \DIFdelbegin \DIFdel{They vary between runs: the floor of a single run is 0.21 to 2.16 times the pooled floor of its geometry, and the alternative-cluster interval is wider than the batch interval. }\DIFdelend Without the warm-up of the first 150 parts the floors fall to 0.75--1.00 of the reference, which keeps the warm-up because series production starts up after every break. Independent part-to-part noise, scanner repeatability included, accounts for 15--33\,\% of each floor, \DIFdelbegin \DIFdel{and the resolution is small against it }\DIFdelend \DIFaddbegin \DIFadd{so A3 holds for repeatability and, }\DIFaddend except for the mid-side draw-in, \DIFdelbegin \DIFdel{so A3 holds for repeatability and }\DIFdelend \DIFaddbegin \DIFadd{for }\DIFaddend resolution, while bias, linearity and stability remain untested.

\DIFdelbegin \DIFdel{Physical candidates for the drift }\DIFdelend \DIFaddbegin \DIFadd{Whether the drift is the process or the scanner cannot be settled without a reference artefact. Its candidates }\DIFaddend are the warm-up of the tools\DIFdelbegin \DIFdel{and its effect on the lubricant, the build-up of lubricant }\DIFdelend \DIFaddbegin \DIFadd{, lubricant building up }\DIFaddend on die and blank holder\DIFdelbegin \DIFdel{over consecutive strokes, and variation of the material and of the blank size }\DIFdelend \DIFaddbegin \DIFadd{, and material and blank size varying }\DIFaddend along the stack\DIFdelbegin \DIFdel{, which enters the draw-in one to one because it is computed against the nominal 210\,mm blank. Whether the drift is the process or the scanner cannot be settled without a reference artefact, but two checks bear on it. No }\DIFdelend \DIFaddbegin \DIFadd{; no }\DIFaddend recorded process signal explains it \DIFdelbegin \DIFdel{: }\DIFdelend \DIFaddbegin \DIFadd{(}\DIFaddend batch centres of \DIFdelbegin \DIFdel{press }\DIFdelend force, punch temperature, \DIFdelbegin \DIFdel{sheet }\DIFdelend thickness and oil film predict it out of fold with $R^2$ of $-0.23$ to $0.20$\DIFdelbegin \DIFdel{. And }\DIFdelend \DIFaddbegin \DIFadd{), and }\DIFaddend the batch centres of the flange width and of the cup width at 15\,mm depth, which punch and die largely fix, are weakly correlated within the runs ($r=-0.27$ and $+0.27$), which \DIFdelbegin \DIFdel{argues against }\DIFdelend \DIFaddbegin \DIFadd{bounds }\DIFaddend a common lateral scale drift of the scanner \DIFdelbegin \DIFdel{; the cup width varies by 0.06--0.07\,mm between batches, so it bounds such a drift only }\DIFdelend below 0.05\,\%, \DIFdelbegin \DIFdel{and }\DIFdelend \DIFaddbegin \DIFadd{about 0.9 floor of mid-side draw-in; }\DIFaddend drift in height and of the edge detection remain untestable. Variation between runs is not in the floor: the three lubrication series of one geometry and force differ by 1.2 (waviness) to 9.8 floors (dome) at the 95th percentile of their centre differences, lubrication effects included, whereas within a run the 95th percentiles of two batches of 100 parts differ by 0.92--1.49 floors in 95\,\% of pairs.

%% BEGIN GENERATED floors
\begin{table}[t]
\tabsize
\setlength{\tabcolsep}{3pt}
\caption{Within-run floor per characteristic, pooled and per geometry, with its spread over runs}\label{tab:floors}
\begin{tabular*}{\textwidth}{@{\extracolsep{\fill}}lrrrrrr@{}}
\toprule
Characteristic & $F$ & Interval & Concave/convex & Single runs & No warm-up & $F/\sigma_\mathrm{LT}$ \\
\midrule
Mid-side (mm) & 0.051 & 0.042--0.061 & 0.050/0.052 & 0.36--1.62 & 0.99 & 1.0/1.4 \\
Corner (mm) & 0.046 & 0.032--0.053 & 0.053/0.030 & 0.27--1.62 & 0.95 & 1.4/1.2 \\
Wavin. (mm) & 0.0028 & 0.0021--0.0050 & 0.0037/0.0021 & 0.21--2.08 & 0.89 & 1.6/0.7 \\
Wall ($^\circ$) & 0.078 & 0.054--0.10 & 0.081/0.078 & 0.30--1.74 & 0.80 & 1.3/0.7 \\
Arm ($^\circ$) & 0.11 & 0.089--0.15 & 0.12/0.100 & 0.40--2.16 & 1.00 & 1.7/0.8 \\
Depth (mm) & 0.011 & 0.0090--0.012 & 0.0087/0.012 & 0.39--1.45 & 0.75 & 1.4/1.6 \\
Dome (mm) & 0.0070 & 0.0052--0.0084 & 0.0034/0.0082 & 0.42--1.61 & 0.93 & 1.2/1.6 \\
\botrule
\end{tabular*}
\footnotetext{Batches of 50 consecutive parts, 95\,\% quantile of the difference of their 10\,\% trimmed means. Interval: 95\,\% over 2000 resamples of whole alternatives (the narrower interval over resampled batches is in the archive). Single runs: floor of one run over that of its geometry, range over the 18 runs. No warm-up: floor without the first 150 parts over the reference. $\sigma_\mathrm{LT}$: standard deviation of the parts of a series (concave/convex); $\sigma_\mathrm{b}$, $\sigma_\mathrm{w}$ and the drift ratio \DIFdelbeginFL \DIFdelFL{are }\DIFdelendFL in \DIFdelbeginFL \DIFdelFL{the archived results}\DIFdelendFL \DIFaddbeginFL \DIFaddFL{Online Resource~1, Table~S4}\DIFaddendFL .}
\end{table}
%% END GENERATED floors

\begin{figure}[t]
\centering
\includegraphics[width=\textwidth]{F_batches}
\caption{The 500 parts of one run in production order with the centres of its ten batches of 50 (bars) and the floor\DIFaddbeginFL \DIFaddFL{, the warm-up shaded}\DIFaddendFL : mid-side draw-in (a) and wall angle (b)}\label{fig:batches}
\end{figure}

\subsection{What production resolves and where the simulation dataset lies}\label{sec:res-sim}

All 63 geometry contrasts exceed the floor and 105 of 126 force contrasts do, but only 51 of 126 lubrication contrasts, so a lubrication variant is often a near-replicate of its siblings. The minimum resolvable force change is 21--67\,kN between 100 and 300\,kN for most characteristics and exceeds either force interval for the wall angle and the dome.

The simulation dataset does not contain the produced process\DIFaddbegin \DIFadd{. The comparisons that depend on neither the assumed blank nor the scanner's scales show it: the simulated force effect on the draw-in is 2.7 to 5.1 times the measured one; the simulated wall angle of the concave cup falls by 1.5$^\circ$ from 100 to 500\,kN, the measured one not at all; the cut arms of the concave cups bend downwards in 89--100\,\% of the simulations but upwards in every part, whereas the convex arm angles agree within 0.6$^\circ$; and the dome ranks the geometries in reverse. The levels, which depend on blank and scales, disagree as well if both hold (\ORt{3})}\DIFaddend : at the nominal thickness the measured centre lies within the range spanned by the friction grid in only 11 of the 42 combinations of characteristic, geometry and force\DIFdelbegin \DIFdel{. The }\DIFdelend \DIFaddbegin \DIFadd{; the }\DIFaddend simulated mid-side draw-in exceeds the measured one at 100\,kN for every friction coefficient (by 1.8--2.0\,mm at friction 0.10) \DIFdelbegin \DIFdel{, reaches }\DIFdelend \DIFaddbegin \DIFadd{and comes within 0.011\,mm of }\DIFaddend it at 300\,kN only at the edge of the grid on the concave cup\DIFdelbegin \DIFdel{, and at 500\,kN at friction 0.105 and 0.116}\DIFdelend ; the corner draw-in exceeds it everywhere, by 2.6--3.5\,mm \DIFdelbegin \DIFdel{. The simulated force effect on the draw-in is 2.7 to 5.1 times the measured one, and the cut arms of the concave cups bend downwards in 89--100}\DIFdelend \DIFaddbegin \DIFadd{at friction 0.10 and 1.3--3.3}\DIFaddend \,\DIFdelbegin \DIFdel{\% of the simulations but upwards in every part, whereas the convex arm angles agree within 0.6$^\circ$}\DIFdelend \DIFaddbegin \DIFadd{mm at 0.15}\DIFaddend . Calibrating friction cannot repair this within the grid: all 72 within-family least-squares calibrations of Mc end at its upper edge (0.15), and with a discrepancy term the posterior of Mcd piles at the lower edge (0.05), where the simulated force effect is weakest. The dataset's mapping from oil film to friction implies 0.8--3.1\,mm of mid-side draw-in per g/m$^2$ of oil, \DIFdelbegin \DIFdel{against within-run slopes of $-0.01$ to $+0.09$\,mm and median slopes }\DIFdelend \DIFaddbegin \DIFadd{an order of magnitude more than the slopes within the runs and }\DIFaddend between the lubrication patterns \DIFdelbegin \DIFdel{of $-0.17$ to $+0.39$}\DIFdelend \DIFaddbegin \DIFadd{(at most 0.09 and 0.39}\DIFaddend \,mm per g/m$^2$\DIFdelbegin \DIFdel{; }\DIFdelend \DIFaddbegin \DIFadd{), even allowing for }\DIFaddend noise in the oil readings\DIFdelbegin \DIFdel{attenuates the within-run slope, but the mapping is too sensitive by an order of magnitude. An }\DIFdelend \DIFaddbegin \DIFadd{; through it an }\DIFaddend envelope over the measured incoming variation \DIFdelbegin \DIFdel{instead of the whole grid halves its width but }\DIFdelend \DIFaddbegin \DIFadd{varies with the lubrication pattern more than production does and }\DIFaddend makes M5 less decisive with a sibling (24 against 7.2 floors)\DIFdelbegin \DIFdel{, because through this mapping it varies with the lubrication pattern more than production does}\DIFdelend . A Gaussian-process surrogate of the simulations decides as the simulation does (109 against 108 floors).

The cup depth is simulated 0.70--1.06\,mm (2.5--3.7\,\%) deeper than measured, and the wall \DIFdelbegin \DIFdel{position }\DIFdelend at 15\,mm depth, set by punch and die, 2.0--2.2\,mm (2.8--3.5\,\%) smaller in radius. The depth offset is \DIFdelbegin \DIFdel{not definitional, since }\DIFdelend \DIFaddbegin \DIFadd{neither definitional, }\DIFaddend bottom and flange \DIFdelbegin \DIFdel{are }\DIFdelend \DIFaddbegin \DIFadd{being }\DIFaddend seen from the same side of the sheet in both, nor due to the flange inclination, which would make the measured depth larger\DIFdelbegin \DIFdel{; a shorter effective drawing depth , a different tool geometry or a scanner scale error could each produce these offsets, which the data cannot separate}\DIFdelend \DIFaddbegin \DIFadd{. A drawing depth 1\,mm shorter would move the wall by only 0.2--0.4\,mm, and a scale error would give an offset proportional to the radius (64 and 75\,mm), not 2.0--2.2\,mm at both; another tool geometry fits best but is unconfirmed}\DIFaddend . Without the \DIFdelbegin \DIFdel{depth }\DIFdelend offset, M0 decides at 65 instead of 108 floors\DIFdelbegin \DIFdel{(a sensitivity only)}\DIFdelend .

\subsection{Decisions by evidence relation}\label{sec:res-decisions}

Table~\ref{tab:distances} and Fig.~\ref{fig:decisions} score the rules; the nominal simulation decides at 108 floors, erring mostly by false rejects.

\emph{A new variant with a produced sibling.} Every calibrated rule improves on it, and the production record is the most decisive: the nearest setting, the mean $q_p$ of the two siblings, is right beyond 3.4 floors (2.0--4.5) and the force trend beyond 3.5, while M5 and M7n decide at 6.8--7.2 floors but are safe below one floor. Stratified by what a team knows when deciding, the nearest setting needs 1.60 floors when the two produced siblings do not differ resolvably \DIFdelbegin \DIFdel{from each other (75 cases on 17 alternatives) }\DIFdelend and 4.40 when they do\DIFdelbegin \DIFdel{(51 cases on 16), the force trend 3.05 and 4.20, and M5 6.15 and 8.90 (strata by the new design's own production, a diagnostic only: 0.85 and 4.85 floors)}\DIFdelend . In physical units the nearest setting is right beyond 0.10\,mm of mid-side draw-in, 0.26$^\circ$ of wall angle and 0.04\,mm of cup depth (Table~\ref{tab:groups}). One produced sibling suffices: with a single alternative at the new design's force the nearest setting decides at 3.5 floors (2.3--4.6), and further alternatives do not improve it, while they make M5 more decisive and safer.

\emph{A new process setting.} With all alternatives at the new force held out, the force trend and the nearest setting decide at 8.2 and 8.5 floors and the interval rules only at 15--19, though they are safe at 5--8. Interpolating is much easier than extrapolating (the nearest setting 4.7 against 9.5 floors), but at 300\,kN, the midpoint of a balanced calibration set, the nearest setting, the force trend and the centre of M5n coincide by construction, so interpolation cannot separate mean-type rules. The scaled simulation, extrapolating a fitted trend, falsely accepts failing alternatives up to 165 floors from their limit, on the cup depth, while on the mid-side draw-in and the arm angle it is more decisive than the force trend \DIFdelbegin \DIFdel{(4.0 against 5.7 and 5.9 against 9.5 floors extrapolating}\DIFdelend \DIFaddbegin \DIFadd{when extrapolating; the multi-fidelity model M7 fails the same way, with false accepts up to 118 floors on the cup depth (\ORt{2}}\DIFaddend ).

\emph{A new family.} Scored in the floor of the produced family, as the procedure prescribes, every rule without the simulation fails (130--135 floors): nothing in one family's record describes the other. With the simulation, Mcd decides at 30 floors and M5 at 35, and M3 is safe at 17 floors and M2 at 26. Per characteristic, M5 decides at 5--10 floors for the draw-ins and the wall angle, which the forming model represents directly, but at 20--40 for depth, waviness, arm and dome. Production knowledge still carries over where a drawing nominal transfers: for the wall angle, the design angle plus the produced family's deviation from its own decides at 1.9 floors, against 7.5 for M5. Two transfers do not resolve which rule is safest.

%% BEGIN GENERATED distances
\begin{table}[t]
\tabsize
\setlength{\tabcolsep}{3.4pt}
\caption{Decisive and safe distances (floors) by the relation of the new design to the produced evidence}\label{tab:distances}
\DIFdelbeginFL %DIFDELCMD < \begin{tabular}{@{}lrrrrrrrr@{}}
%DIFDELCMD < \toprule
%DIFDELCMD <  & \multicolumn{2}{c}{Sibling} & \multicolumn{2}{c}{Interpolated} & \multicolumn{2}{c}{Extrapolated} & \multicolumn{2}{c}{New family} \\
%DIFDELCMD < \cmidrule(lr){2-3}\cmidrule(lr){4-5}\cmidrule(lr){6-7}\cmidrule(lr){8-9}
%DIFDELCMD < Rule & $\delta_\mathrm{d}/\delta_\mathrm{s}$ & FA/FR & $\delta_\mathrm{d}/\delta_\mathrm{s}$ & FA/FR & $\delta_\mathrm{d}/\delta_\mathrm{s}$ & FA/FR & $\delta_\mathrm{d}/\delta_\mathrm{s}$ & FA/FR \\
%DIFDELCMD < \midrule
%DIFDELCMD < M0 & 108 & 17/109 & 109 & 17/115 & 106 & 12/109 & 104 & 28/105 \\
%DIFDELCMD < M1s$^\dagger$ & 4.9 & 4.7/4.9 & 19 & 30/19 & 19 & 165/7.9 & 121 & 125/118 \\
%DIFDELCMD < M2 & 23/$<$1 & $<$1/$<$1 & 18/2.7 & 3.7/2.4 & 18/9.9 & 8.7/10.0 & 41/26 & 27/9.7 \\
%DIFDELCMD < M5 & 7.2/$<$1 & $<$1/$<$1 & 19/2.8 & 2.4/3.1 & 19/8.9 & 8.4/9.7 & 35/31 & 33/15 \\
%DIFDELCMD < \quad interval of $\delta_\mathrm{d}$ & 6.6--8.4 &  & 18--20 &  & 18--20 &  & 31--39 & \\
%DIFDELCMD < M7n$^\ddagger$ & 6.8/$<$1 & $<$1/$<$1 & 7.4/3.4 & 2.2/3.4 & 15/5.9 & 7.5/$<$1 & 132/130 & 129/167 \\
%DIFDELCMD < \quad interval of $\delta_\mathrm{d}$ & 6.0--7.3 &  & 7.1--8.1 &  & 15--16 &  & 131--132 & \\
%DIFDELCMD < NN$^\dagger$ & 3.4 & 2.8/3.4 & 4.7 & 3.1/4.7 & 9.5 & 11/5.6 & 130 & 129/168 \\
%DIFDELCMD < \quad interval of $\delta_\mathrm{d}$ & 2.0--4.5 &  & 4.3--5.1 &  & 6.6--11 &  & 129--131 & \\
%DIFDELCMD < M1sn$^\dagger$ & 3.5 & 3.5/3.6 & 4.7 & 3.1/4.7 & 8.9 & 9.3/7.4 & 131 & 129/168 \\
%DIFDELCMD < NN+RM$^\ddagger$ & 6.3/$<$1 & $<$1/$<$1 & 9.6/$<$1 & $<$1/$<$1 & 20/$<$1 & 1.5/$<$1 & 293/4.5 & 4.1/44 \\
%DIFDELCMD < \midrule Trial, first 20 parts & 1.4 & 1.5/1.2 & 1.4 & 1.5/1.2 & 1.4 & 1.5/1.2 & 1.4 & 1.5/1.2 \\
%DIFDELCMD < Trial, first 50 parts & 0.90 & 1.0/0.85 & 0.90 & 1.0/0.85 & 0.90 & 1.0/0.85 & 0.90 & 1.0/0.85 \\
%DIFDELCMD < \botrule
%DIFDELCMD < \end{tabular}
%DIFDELCMD < %%%
\DIFdelendFL \DIFaddbeginFL \begin{tabular}{@{}lrrrrrrrr@{}}
\toprule
 & \multicolumn{2}{c}{Sibling} & \multicolumn{2}{c}{Interpolated} & \multicolumn{2}{c}{Extrapolated} & \multicolumn{2}{c}{New family} \\
\cmidrule(lr){2-3}\cmidrule(lr){4-5}\cmidrule(lr){6-7}\cmidrule(lr){8-9}
Rule & $\delta_\mathrm{d}/\delta_\mathrm{s}$ & FA/FR & $\delta_\mathrm{d}/\delta_\mathrm{s}$ & FA/FR & $\delta_\mathrm{d}/\delta_\mathrm{s}$ & FA/FR & $\delta_\mathrm{d}/\delta_\mathrm{s}$ & FA/FR \\
\midrule
M0 & 108 & 17/109 & 109 & 17/115 & 106 & 12/109 & 104 & 28/105 \\
M1s$^\dagger$ & 4.9 & 4.7/4.9 & 19 & 30/19 & 19 & 165/7.9 & 121 & 125/118 \\
M2 & 23/$<$1 & $<$1/$<$1 & 18/2.7 & 3.7/2.4 & 18/9.9 & 8.7/10.0 & 41/26 & 27/9.7 \\
M5 & 7.2/$<$1 & $<$1/$<$1 & 19/2.8 & 2.4/3.1 & 19/8.9 & 8.4/9.7 & 35/31 & 33/15 \\
\quad interval of $\delta_\mathrm{d}$ & 6.6--8.4 &  & 18--20 &  & 18--20 &  & 31--39 & \\
M7n$^\ddagger$ & 6.8/$<$1 & $<$1/$<$1 & 7.4/3.4 & 2.2/3.4 & 15/5.9 & 7.5/$<$1 & 132/130 & 129/167 \\
\quad interval of $\delta_\mathrm{d}$ & 6.0--7.3 &  & 7.1--8.1 &  & 15--16 &  & 131--132 & \\
NN$^\dagger$ & 3.4 & 2.8/3.4 & 4.7 & 3.1/4.7 & 9.5 & 11/5.6 & 130 & 129/168 \\
\quad interval of $\delta_\mathrm{d}$ & 2.0--4.5 &  & 4.3--5.1 &  & 6.6--11 &  & 129--131 & \\
M1sn$^\dagger$ & 3.5 & 3.5/3.6 & 4.7 & 3.1/4.7 & 8.9 & 9.3/7.4 & 131 & 129/168 \\
NN+RM$^\ddagger$ & 6.3/$<$1 & $<$1/$<$1 & 9.6/$<$1 & $<$1/$<$1 & 20/$<$1 & 1.5/$<$1 & 293/4.5 & 4.1/44 \\
Mcd$^\ddagger$ & 7.2/$<$1 & $<$1/$<$1 & 22/1.6 & $<$1/2.5 & 24/5.8 & 6.6/2.6 & 30/19 & 29/12 \\
\midrule Trial, first 20 parts & 1.4 & 1.5/1.2 & 1.4 & 1.5/1.2 & 1.4 & 1.5/1.2 & 1.4 & 1.5/1.2 \\
Trial, first 50 parts & 0.90 & 1.0/0.85 & 0.90 & 1.0/0.85 & 0.90 & 1.0/0.85 & 0.90 & 1.0/0.85 \\
\botrule
\end{tabular}
\DIFaddendFL \footnotetext{Over all seven characteristics; one number for rules that always decide; $<$1: below one floor. FA/FR: safe distances on the false-accept and the false-reject side. 126 cases with a sibling and for a new family, 42 interpolated, 84 extrapolated; a new family in the produced family's floor (in its own floor\DIFaddbeginFL \DIFaddFL{, and all rules, }\DIFaddendFL in \DIFdelbeginFL \DIFdelFL{the archived results}\DIFdelendFL \DIFaddbeginFL \DIFaddFL{Online Resource~1, Table~S1}\DIFaddendFL ). Interval: 95\,\% over 1000 resamples of the held-out alternatives, fits fixed. \DIFaddbeginFL \DIFaddFL{Mcd: friction calibration with a discrepancy term (Section~\ref{sec:rules}). }\DIFaddendFL Trial: the first $n$ parts of the design's own run, a within-run lower bound. $^\dagger$, $^\ddagger$: added after the first results and \DIFaddbeginFL \DIFaddFL{in the revisions }\DIFaddendFL for this \DIFdelbeginFL \DIFdelFL{version}\DIFdelendFL \DIFaddbeginFL \DIFaddFL{journal}\DIFaddendFL .}
\end{table}
%% END GENERATED distances

%% BEGIN GENERATED groups
\begin{table}[t]
\tabsize
\setlength{\tabcolsep}{1pt}
\caption{Decisive distance per characteristic in floors and, in brackets, in its unit}\label{tab:groups}
\begin{tabular*}{\textwidth}{@{\extracolsep{\fill}}lrrrrrrr@{}}
\toprule
 & \multicolumn{3}{c}{NN} & M1sn & M1s & M5 & \\
\cmidrule(lr){2-4}
Characteristic & Sibling & Interp. & Extrap. & Extrap. & Extrap. & New family & $P_\mathrm{pk}$ \\
\midrule
Mid-side (mm) & 2.0 (0.099) & 3.8 (0.19) & 6.5 (0.32) & 5.7 (0.28) & 4.0 (0.20) & 5.5 (0.27) & 1.3 \\
Corner (mm) & 4.5 (0.13) & 1.9 (0.058) & 6.7 (0.23) & 2.8 (0.083) & 2.9 (0.100) & 9.9 (0.31) & 2.5 \\
Wavin. (mm) & $<$1 ($<$0.0028) & 4.8 (0.014) & 11 (0.040) & 9.3 (0.027) & 21 (0.077) & 22 (0.047) & 0.9 \\
Wall ($^\circ$) & 3.4 (0.26) & $<$1 ($<$0.078) & 1.9 (0.15) & 1.8 (0.14) & 2.0 (0.16) & 7.5 (0.58) & 1.7 \\
Arm ($^\circ$) & 2.8 (0.28) & 5.9 (0.59) & 5.4 (0.53) & 9.5 (0.94) & 5.9 (0.69) & 39 (3.9) & 1.6 \\
Depth (mm) & 3.3 (0.039) & 4.4 (0.038) & 15 (0.13) & 9.9 (0.086) & 165 (1.4) & 20 (0.20) & 2.4 \\
Dome (mm) & 7.0 (0.038) & 5.3 (0.044) & 4.3 (0.020) & 3.0 (0.020) & 15 (0.056) & 40 (0.21) & 4.0 \\
\botrule
\end{tabular*}
\footnotetext{18 cases per characteristic with a sibling and for a new family, 6 interpolated, 12 extrapolated; \DIFdelbeginFL \DIFdelFL{units through the pooled floor of the characteristic; }\DIFdelendFL a new family in the produced family's floor. \DIFaddbeginFL \DIFaddFL{In brackets: the decisive distance computed in the unit, every case scored in the pooled floor (not the distance in floors times a floor). }\DIFaddendFL $P_\mathrm{pk}$: median performance index of a requirement at the decisive distance of NN with a sibling. All \DIFaddbeginFL \DIFaddFL{main }\DIFaddendFL rules \DIFdelbeginFL \DIFdelFL{and the safe }\DIFdelendFL \DIFaddbeginFL \DIFaddFL{with both }\DIFaddendFL distances \DIFdelbeginFL \DIFdelFL{are }\DIFdelendFL in \DIFdelbeginFL \DIFdelFL{the archived results}\DIFdelendFL \DIFaddbeginFL \DIFaddFL{Online Resource~1, Table~S2}\DIFaddendFL .}
\end{table}
%% END GENERATED groups

\begin{figure}[t]
\centering
\includegraphics[width=\textwidth]{F3_decisions_main}
\caption{Verdict shares against the requirement distance $d$ (symmetric logarithmic axis, linear within $\pm10$ floors; 126 cases per panel; a new family in the produced family's floor)}\label{fig:decisions}
\end{figure}

\subsection{Learning given a fair test}\label{sec:res-ml}

The Gaussian processes fitted by maximum marginal likelihood reduce to the baselines in this design \DIFaddbegin \DIFadd{(\ORt{5})}\DIFaddend . For a new setting the force length scale \DIFaddbegin \DIFadd{and noise variance }\DIFaddend of M5n \DIFdelbegin \DIFdel{sits at its lower bound in 68--79\,\% of the fits and its noise variance at its bound in every fit, }\DIFdelend \DIFaddbegin \DIFadd{sit at their bounds in most fits }\DIFaddend and its centre lies within 0.01 floor of the calibration mean in 93--95\,\% of the cases\DIFdelbegin \DIFdel{; }\DIFdelend \DIFaddbegin \DIFadd{, }\DIFaddend with a sibling \DIFdelbegin \DIFdel{it lies }\DIFdelend a median 0.28 floor from the nearest setting. The trend GP, specified for this regime, reduces to the force trend instead: its centre lies a median 0.33--0.36 floor from it for a new setting, its force length scale stays near the prior's 200\,kN, which two force levels cannot inform, and its between-run nugget, 0.4--0.5 floors, is three to four times the sampling error of the truths. It is more decisive than M5n for a new setting (14.5 against 17.2 floors\DIFdelbegin \DIFdel{; paired difference $-2.7$, $-5.6$ to $-1.1$}\DIFdelend ) and as safe, but less decisive than the nearest setting in every relation (by 3.5 floors with a sibling and 6 for a new setting). \DIFaddbegin \DIFadd{Its intervals cover only 52\,\% interpolating and 56\,\% extrapolating (17 and 33\,\% on the mid-side draw-in, none on the waviness), because their width is learned from the variation between lubrication patterns at produced forces, sibling-type information; halving or doubling any of its four prior medians moves its decisive distances by at most 1.4 floors and keeps it within 0.4 floor of the force trend and its coverage for a new setting at 48--62\,\% (\ORt{6}). }\DIFaddend With the simulation in its mean (M7) it gets worse for a new setting (28 floors), because the simulated force trend enters it. \DIFdelbegin \DIFdel{The archive's study of }\DIFdelend \DIFaddbegin \DIFadd{Tuned }\DIFaddend part-level boosting \DIFdelbegin \DIFdel{(regularised, flexible, ridge, quantile loss, inner selection) finds no setting }\DIFdelend \DIFaddbegin \DIFadd{decides no }\DIFaddend closer than 11.2 floors with a sibling and 13.1 for a new setting.

Prediction and interval separate. A relation-matched interval borrowed from the other family makes even the nearest setting safe for a new setting (safe distance $<$1 floor, against 6.8 for M5 and 5.2 for M7n) at the decisive distance of the learned intervals (19.5 against 18.7 and 14.5 floors)\DIFdelbegin \DIFdel{, }\DIFdelend \DIFaddbegin \DIFadd{. The safety comes from the borrowed evidence, not from the predictor (\ORt{7}): around the trend GP and force-trend centres the same interval decides closer when extrapolating (17.6 }\DIFaddend and \DIFdelbegin \DIFdel{covers }\DIFdelend \DIFaddbegin \DIFadd{17.8 against 20.1 floors), around simulation-based centres it does not make the rule safe (false accepts up to 151--155 floors), and a trend GP calibrated on both families, borrowing through a family term, is safe at 2.75 floors. The coverage of NN+RM, }\DIFaddend 88\,\% \DIFdelbegin \DIFdel{of the truths }\DIFdelend interpolating and 93\,\% extrapolating, \DIFdelbegin \DIFdel{against 71 and 46}\DIFdelend \DIFaddbegin \DIFadd{comes from two family swaps, not from a guarantee: it ranges from 73 to 100}\DIFaddend \,\% \DIFdelbegin \DIFdel{for M5 and 52 and 56}\DIFdelend \DIFaddbegin \DIFadd{between the families and from 61 to 100}\DIFaddend \,\% \DIFdelbegin \DIFdel{for M7n. By the interval score \mbox{%DIFAUXCMD
\citep{gneiting2007strictly} }\hskip0pt%DIFAUXCMD
}\DIFdelend \DIFaddbegin \DIFadd{between characteristics (\ORt{9}). Around production-record and trend centres }\DIFaddend the relation-matched intervals \DIFdelbegin \DIFdel{are the best }\DIFdelend \DIFaddbegin \DIFadd{have the best interval scores \mbox{%DIFAUXCMD
\citep{gneiting2007strictly} }\hskip0pt%DIFAUXCMD
}\DIFaddend for a new setting (\DIFdelbegin \DIFdel{12--14 interpolating and 19--24 extrapolating, against 14--98 and 41--147 for the rules' own intervals}\DIFdelend \DIFaddbegin \DIFadd{\ORt{8}}\DIFaddend ). With a sibling the \DIFdelbegin \DIFdel{model-based intervals (M5, M5n, M7n, M7, Mcd) }\DIFdelend \DIFaddbegin \DIFadd{Gaussian-process intervals }\DIFaddend cover 83--86\,\% (cluster intervals \DIFdelbegin \DIFdel{0.71}\DIFdelend \DIFaddbegin \DIFadd{0.72}\DIFaddend --0.94), compatible with their nominal 90\,\%, and the jackknife+ intervals of M3 96\,\% (0.93--0.98) against a guarantee of 0.78\DIFdelbegin \DIFdel{; the conformal margin of M6, computed on sibling-type residuals, covers 93\,\% with a sibling at a cost of ten floors of decisiveness against M5 (17 against 7.2) but 61\,\% extrapolating}\DIFdelend . For a new family no interval covers well except at widths that send every design to trial\DIFdelbegin \DIFdel{. With one produced alternative }\DIFdelend \DIFaddbegin \DIFadd{; with one to three produced alternatives }\DIFaddend of the new family, its \DIFaddbegin \DIFadd{own nearest setting or its }\DIFaddend offset to the other family's record decides at \DIFdelbegin \DIFdel{8.9 }\DIFdelend \DIFaddbegin \DIFadd{5.9--10.5 }\DIFaddend floors in the \DIFdelbegin \DIFdel{new family's floor and the family's own record at 10.5; with two or three, its own nearest setting decides at 7.7 and 5.9 floors, while a two-task Gaussian process with a free inter-family correlation is safe at 1.3--1.6 floors but decides only at 13--17}\DIFdelend \DIFaddbegin \DIFadd{floor of its nine runs, which a team with so few alternatives does not have}\DIFaddend .

%% BEGIN GENERATED guards
\begin{table}[t]
\tabsize
\setlength{\tabcolsep}{2.2pt}
\caption{Guard bands \DIFdelbeginFL \DIFdelFL{of step 6 }\DIFdelendFL \DIFaddbeginFL \DIFaddFL{available when deciding }\DIFaddendFL (floors), \DIFaddbeginFL \DIFaddFL{with }\DIFaddendFL their \DIFdelbeginFL \DIFdelFL{out-of-sample }\DIFdelendFL share correct and \DIFdelbeginFL \DIFdelFL{the }\DIFdelendFL trial rate}\label{tab:guards}
\DIFdelbeginFL %DIFDELCMD < \begin{tabular*}{\textwidth}{@{\extracolsep{\fill}}lrrrrrrrrrrrr@{}}
%DIFDELCMD < \toprule
%DIFDELCMD <  & \multicolumn{4}{c}{Sibling} & \multicolumn{4}{c}{Interpolated} & \multicolumn{4}{c}{Extrapolated} \\
%DIFDELCMD < \cmidrule(lr){2-5}\cmidrule(lr){6-9}\cmidrule(lr){10-13}
%DIFDELCMD < Rule & $g$ & $g_5$ & OOS & T & $g$ & $g_5$ & OOS & T & $g$ & $g_5$ & OOS & T \\
%DIFDELCMD < \midrule
%DIFDELCMD < NN$^\dagger$ & 0.25 & 2.0 & 95 & 2 & 2.0 & 4.25 & 95 & 18 & 7.0 & 14.25 & 94 & 55 \\
%DIFDELCMD < M1sn$^\dagger$ & 0.75 & 2.75 & 95 & 7 & 2.0 & 4.25 & 95 & 18 & 6.0 & 8.5 & 95 & 48 \\
%DIFDELCMD < M2 & 0 & 0 & 99 & 54 & 0.75 & 2.75 & 95 & 53 & none & none & -- & 100 \\
%DIFDELCMD < M5 & 0 & 0 & 99 & 23 & 0.5 & 2.75 & 96 & 54 & none & none & -- & 100 \\
%DIFDELCMD < M7n$^\ddagger$ & 0 & 6.75 & 98 & 19 & 0.75 & 3.75 & 95 & 27 & 3.25 & 6.75 & 95 & 56 \\
%DIFDELCMD < NN+RM$^\ddagger$ & 0 & 0.5 & 99 & 28 & 0 & 0 & 99 & 42 & 0 & 5.25 & 97 & 72 \\
%DIFDELCMD < \botrule
%DIFDELCMD < \end{tabular*}
%DIFDELCMD < %%%
\DIFdelendFL \DIFaddbeginFL \begin{tabular*}{\textwidth}{@{\extracolsep{\fill}}lrrrrrrrrr@{}}
\toprule
 & \multicolumn{3}{c}{Sibling} & \multicolumn{3}{c}{Interpolated} & \multicolumn{3}{c}{Extrapolated} \\
\cmidrule(lr){2-4}\cmidrule(lr){5-7}\cmidrule(lr){8-10}
Rule & $g$ & Correct & T & $g$ & Correct & T & $g$ & Correct & T \\
\midrule
\multicolumn{10}{@{}l}{\emph{Reference prior}} \\
NN$^\dagger$ & 0/0.5 & 95 & 2 & 2.5/1.5 & 95 & 18 & 5.75/9.5 & 95 & 58 \\
M1sn$^\dagger$ & 0.5/1.0 & 95 & 7 & 2.5/1.5 & 95 & 18 & 6.5/5.5 & 95 & 47 \\
M2 & 0 & 99 & 54 & 0.5/0.75 & 95 & 52 & 6.25/none & 84 & 81 \\
M5 & 0 & 99 & 23 & 0/1.75 & 95 & 57 & 5.0/none & 83 & 80 \\
M7n$^\ddagger$ & 0 & 98 & 19 & 1.75/0 & 95 & 28 & 4.75/2.75 & 96 & 60 \\
NN+RM$^\ddagger$ & 0 & 99 & 28 & 0 & 99 & 42 & 0/0.5 & 97 & 73 \\
\multicolumn{10}{@{}l}{\emph{Release-like prior ($\pm5$ floors)}} \\
NN$^\dagger$ & 1.0/3.0 & 96 & 41 & 4.75/3.5 & 93 & 67 & 8.5/14.5 & 75 & 92 \\
M1sn$^\dagger$ & 2.25/3.0 & 95 & 47 & 4.75/3.5 & 93 & 67 & 8.5/8.25 & 94 & 84 \\
M2 & 0 & 97 & 76 & 3.25/2.25 & 94 & 85 & 10.25/none & 72 & 91 \\
M5 & 0/6.0 & 95 & 67 & 1.5/3.25 & 88 & 84 & 9.25/none & 67 & 91 \\
M7n$^\ddagger$ & 5.5/7.25 & 84 & 91 & 4.25/1.75 & 87 & 74 & 7.0/6.25 & 93 & 90 \\
NN+RM$^\ddagger$ & 0/1.0 & 96 & 72 & 0.5/0 & 95 & 77 & 0/5.25 & 75 & 93 \\
\botrule
\end{tabular*}
\DIFaddendFL \footnotetext{$g$\DIFdelbeginFL \DIFdelFL{, $g_5$}\DIFdelendFL : \DIFdelbeginFL \DIFdelFL{guard }\DIFdelendFL band \DIFdelbeginFL \DIFdelFL{under }\DIFdelendFL \DIFaddbeginFL \DIFaddFL{of }\DIFaddendFL the \DIFdelbeginFL \DIFdelFL{reference prior and under }\DIFdelendFL \DIFaddbeginFL \DIFaddFL{family's own sibling cases with a sibling; for a new setting, which a family produced at two forces cannot exercise, }\DIFaddendFL the \DIFdelbeginFL \DIFdelFL{release-like prior }\DIFdelendFL \DIFaddbeginFL \DIFaddFL{band of the same relation in the other family }\DIFaddendFL (\DIFdelbeginFL \DIFdelFL{requirements within $\pm5$ floors)}\DIFdelendFL \DIFaddbeginFL \DIFaddFL{concave/convex; one number when equal}\DIFaddendFL ; none: no band up to 20 floors reaches 95\,\%\DIFaddbeginFL \DIFaddFL{)}\DIFaddendFL . \DIFdelbeginFL \DIFdelFL{OOS}\DIFdelendFL \DIFaddbeginFL \DIFaddFL{Correct}\DIFaddendFL : share correct \DIFdelbeginFL \DIFdelFL{(\%) }\DIFdelendFL among decided verdicts \DIFaddbeginFL \DIFaddFL{(\%), }\DIFaddendFL with \DIFdelbeginFL \DIFdelFL{the band estimated }\DIFdelendFL \DIFaddbeginFL \DIFaddFL{a sibling }\DIFaddendFL leaving \DIFdelbeginFL \DIFdelFL{each }\DIFdelendFL \DIFaddbeginFL \DIFaddFL{one }\DIFaddendFL alternative out within its family\DIFaddbeginFL \DIFaddFL{, for a new setting the other family's band applied to this family}\DIFaddendFL . T: share of verdicts sent to trial (\%)\DIFdelbeginFL \DIFdelFL{, reference prior}\DIFdelendFL . No rule qualifies for a new family\DIFdelbeginFL \DIFdelFL{; bands per family}\DIFdelendFL \DIFaddbeginFL \DIFaddFL{. Intervals}\DIFaddendFL , \DIFdelbeginFL \DIFdelFL{out of sample }\DIFdelendFL \DIFaddbeginFL \DIFaddFL{verdict counts }\DIFaddendFL and \DIFdelbeginFL \DIFdelFL{under other priors are in }\DIFdelendFL the \DIFdelbeginFL \DIFdelFL{archived results}\DIFdelendFL \DIFaddbeginFL \DIFaddFL{characteristics in Online Resource~1, Tables~S10 and S11}\DIFaddendFL .}
\end{table}
%% END GENERATED guards

\subsection{Guard bands}\label{sec:res-procedure}

Step 5 turns the margins into guard bands \DIFaddbegin \DIFadd{as a team could estimate them when deciding }\DIFaddend (Table~\ref{tab:guards})\DIFaddbegin \DIFadd{: with a sibling from the family's own sibling cases, and for a new setting from the same relation in the other family, since a family produced at two forces has no interpolation case of its own}\DIFaddend . Under the reference prior the nearest setting \DIFdelbegin \DIFdel{needs 0.25 floors with a sibling, 2.0 interpolating and 7.0 extrapolating, }\DIFdelend \DIFaddbegin \DIFadd{keeps 95\,\% of its decided verdicts correct in each relation }\DIFaddend and sends 2, 18 and \DIFdelbegin \DIFdel{55}\DIFdelend \DIFaddbegin \DIFadd{58}\DIFaddend \,\% of the verdicts to trial; M2 and M5 \DIFdelbegin \DIFdel{need no band with a sibling and less than one floor interpolating, but send 23--54}\DIFdelend \DIFaddbegin \DIFadd{send 52--57}\DIFaddend \,\% \DIFdelbegin \DIFdel{of the verdicts to trial , and they do not qualify for an extrapolated setting. Estimated per family, as step 5 prescribes, the bands differ between the geometries (the nearest setting 1.5 }\DIFdelend \DIFaddbegin \DIFadd{to trial interpolating }\DIFaddend and\DIFdelbegin \DIFdel{2.5 floors interpolating }\DIFdelend \DIFaddbegin \DIFadd{, extrapolating}\DIFaddend , \DIFdelbegin \DIFdel{9.5 and5.75 extrapolating), and M2 and M5 qualify for an extrapolated setting }\DIFdelend \DIFaddbegin \DIFadd{qualify only }\DIFaddend on the convex cup\DIFdelbegin \DIFdel{only; moved to the concave cup, their bands give only 84 and 83}\DIFdelend \DIFaddbegin \DIFadd{, whose band gives 83--84}\DIFaddend \,\% correct \DIFdelbegin \DIFdel{decided verdicts. Out of sample the bands hold: leaving one alternative out keeps the share of correct decided verdicts at 0.94--0.99 for every rule and relation, and leaving one force cell out at 0.93--0.95 for the nearest setting and the force trend, but at 0.83 and 0.81 for M2 and M5 when extrapolating}\DIFdelend \DIFaddbegin \DIFadd{on the concave cup}\DIFaddend . The 95\,\% is an average over characteristics: \DIFdelbegin \DIFdel{with a sibling }\DIFdelend the nearest setting is right in 89\,\% of its decided verdicts on the dome \DIFdelbegin \DIFdel{, 98.5}\DIFdelend \DIFaddbegin \DIFadd{with a sibling, 88}\DIFaddend \,\% on the waviness \DIFdelbegin \DIFdel{and 92}\DIFdelend \DIFaddbegin \DIFadd{interpolating and 85}\DIFaddend \,\% on the \DIFdelbegin \DIFdel{cut part}\DIFdelend \DIFaddbegin \DIFadd{cup depth extrapolating (\ORt{11}). The bands are within-run: with a between-run component of $\tau=3$ floors the pooled bands of the nearest setting grow from 0.25, 2.0 and 7.0 to 2.9, 4.0 and 8.9 floors (\ORt{14})}\DIFaddend .

The bands belong to the prior as much as to the rules. Under the release-like prior \DIFdelbegin \DIFdel{the band of }\DIFdelend \DIFaddbegin \DIFadd{they widen, and }\DIFaddend the nearest setting \DIFdelbegin \DIFdel{grows to 2.0 floors with a sibling, 4.25 interpolating and 14.25 extrapolating, and it decides only 63, 32 and 2}\DIFdelend \DIFaddbegin \DIFadd{sends 41, 67 and 92\,\% of the verdicts to trial and keeps only 75}\DIFaddend \,\% of \DIFdelbegin \DIFdel{the verdicts ; leaving a force cellout, its share correct drops to }\DIFdelend \DIFaddbegin \DIFadd{its decided verdicts correct extrapolating; with the family's own band estimated without the held-out force cell, the nearest setting and NN+RM reach }\DIFaddend 0.72 \DIFdelbegin \DIFdel{when extrapolating}\DIFdelend \DIFaddbegin \DIFadd{and 0.74 extrapolating, M5 and M7n 0.88 and 0.87 interpolating (\ORt{10})}\DIFaddend . Under the reference prior no rule qualifies for a new family, as the practice of trying out every new die would have it: Mcd reaches 95\,\% on one transfer with a band of 18 floors, but that band gives 73\,\% on the other\DIFdelbegin \DIFdel{(under a uniform prior within $\pm50$ floors M2, M3 and M5 qualify). The applicability guards, reported as diagnostics, refuse indiscriminately: the range guard refuses exactly the extrapolated designs, and 96\,\% of its refusals at requirements 10 floors from the truth would have been right for the nearest setting}\DIFdelend .

\subsection{Costs, trials and release-level requirements}\label{sec:res-costs}

\DIFdelbegin \DIFdel{Pricing }\DIFdelend \DIFaddbegin \DIFadd{Fig.~\ref{fig:costs} prices }\DIFaddend the verdicts by Eq.~(\ref{eq:cost})\DIFdelbegin \DIFdel{with }\DIFdelend \DIFaddbegin \DIFadd{, }\DIFaddend a false reject \DIFdelbegin \DIFdel{at half the cost of }\DIFdelend \DIFaddbegin \DIFadd{costing half }\DIFaddend a false accept\DIFdelbegin \DIFdel{(Fig.~\ref{fig:costs}), the production record and the Gaussian processes cost about the same with a sibling while trials are cheap, and }\DIFdelend \DIFaddbegin \DIFadd{; the following holds under both priors and every ratio of false-reject to false-accept cost (\ORt{13}). For a new setting the nearest setting or the force trend is the cheapest rule at every trial cost. With a sibling }\DIFaddend the nearest setting is the cheapest \DIFdelbegin \DIFdel{rule }\DIFdelend once a trial costs a \DIFdelbegin \DIFdel{tenth }\DIFdelend \DIFaddbegin \DIFadd{fifth }\DIFaddend of a false accept or more; \DIFdelbegin \DIFdel{for a new setting the force trend or the nearest setting is the cheapest at every trial cost; and for }\DIFdelend \DIFaddbegin \DIFadd{below that the cheapest rules differ by thousandths of a false accept, which this design does not resolve. For }\DIFaddend a new family trying every design is the cheapest \DIFdelbegin \DIFdel{until a trial costs about a tenth }\DIFdelend \DIFaddbegin \DIFadd{up to a twentieth }\DIFaddend of a false accept, beyond which \DIFdelbegin \DIFdel{M2 and then M5 are. The }\DIFdelend \DIFaddbegin \DIFadd{the cheapest rule changes with the prior and the cost ratio; minimised without the 95\,\% target, the cost even favours the nominal simulation, with false accepts in 14\,\% of its verdicts. The }\DIFaddend cost-minimising \DIFdelbegin \DIFdel{band of the nearest setting falls from 5 to 0 floors with a sibling and from 15 to 0 when extrapolating as the trial cost rises from 0.01 to 0.5 of a false accept; }\DIFdelend \DIFaddbegin \DIFadd{bands, in sample and pooled over the families, widen as trials get cheaper, and }\DIFaddend the 95\,\% bands of step 5 correspond to trials costing a tenth to a quarter of a false accept\DIFdelbegin \DIFdel{, more than }\DIFdelend \DIFaddbegin \DIFadd{; if, as assumed here, trying }\DIFaddend a new force on an existing tool costs \DIFdelbegin \DIFdel{, so }\DIFdelend \DIFaddbegin \DIFadd{less, }\DIFaddend cost-optimal bands \DIFdelbegin \DIFdel{would be }\DIFdelend \DIFaddbegin \DIFadd{are }\DIFaddend wider. A trial\DIFdelbegin \DIFdel{is not }\DIFdelend \DIFaddbegin \DIFadd{, priced as }\DIFaddend error-free \DIFdelbegin \DIFdel{either}\DIFdelend \DIFaddbegin \DIFadd{in Fig.~\ref{fig:costs} (whose band search runs to 40 floors), is not}\DIFaddend : scored as a rule, a trial of the first 20 or 50 parts of the design's own run decides at 1.35 and 0.90 floors, \DIFaddbegin \DIFadd{and the errors of the 20-part trial cost 0.017 of a false accept per verdict (0.037 under the release-like prior), as much as the whole expected cost of the nearest setting with a sibling at }\DIFaddend a \DIFdelbegin \DIFdel{lower bound, since a trial on another day also carries the variation between runs.
}\DIFdelend \DIFaddbegin \DIFadd{trial cost of 0.05 (0.018). Both priors place 47\,\% of the verdicts on failing designs.
}\DIFaddend 

\DIFdelbegin \DIFdel{Upper limits at which an alternative has a performance index of 1.0, 1.33, 1.67 or 2.0 lie a median 1.03, 1.75, 2.46 and 3.14 floors above its 95th percentile. }\DIFdelend A limit at \DIFdelbegin \DIFdel{an }\DIFdelend \DIFaddbegin \DIFadd{a performance }\DIFaddend index of 1.33, \DIFdelbegin \DIFdel{a recommended minimum for existing processes and }\DIFdelend the lower end of the range in which an initial process study may be acceptable after a customer review, \DIFdelbegin \DIFdel{therefore lies }\DIFdelend \DIFaddbegin \DIFadd{lies a median 1.75 floors above an alternative's 95th percentile (1.03, 2.46 and 3.14 at indices of 1.0, 1.67 and 2.0), }\DIFaddend inside every attainable distance except \DIFdelbegin \DIFdel{that of }\DIFdelend the nearest setting\DIFaddbegin \DIFadd{'s }\DIFaddend when the produced siblings do not differ resolvably \DIFdelbegin \DIFdel{from each other }\DIFdelend (1.60 floors)\DIFaddbegin \DIFadd{. Release, however, asks whether the index is at least 1.33, a decision on $\bar m+4\sigma$: scored on it, the nearest setting needs 3.35 floors with a sibling and 10.55 extrapolating, so it confirms an index of 1.33 only for designs whose true within-run index is 2.1--3.3 with a sibling and 3.8--7.4 extrapolating. Nor does a short trial settle the question once runs differ: with a between-run component of 1}\DIFaddend , \DIFdelbegin \DIFdel{which itexceeds by only 0.15 floor: release-level requirements are decided at the design stage only with a near-replicate produced sibling , if at all}\DIFdelend \DIFaddbegin \DIFadd{2 and 3 floors, the 20-part trial decides at 2.1, 3.6 and 5.1 floors, the nearest setting with a sibling at 3.45, 4.45 and 5.75 (\ORt{14}). Near release level the evidence must come from replicated runs}\DIFaddend .

\begin{figure}[t]
\centering
\includegraphics[width=\textwidth]{F_costs}
\caption{Expected cost per verdict with the cost-minimising band (a--d) and that band (e--h) against the cost of a trial; reference prior, $c_\mathrm{FR}=0.5\,c_\mathrm{FA}$\DIFaddbeginFL \DIFaddFL{, error-free trials, bands up to 40 floors}\DIFaddendFL }\label{fig:costs}
\end{figure}

\subsection{Sensitivity, selection and robustness}\label{sec:res-robust}

A between-run component added to every truth lengthens the distances but keeps their order: with $\tau=3$ floors, about the largest variation between series, the nearest setting needs 5.8 floors with a sibling and 10 extrapolating, M5 9.3 and 20, \DIFaddbegin \DIFadd{and }\DIFaddend a production-record rule (nearest setting, force trend or median) is the most decisive within the family in each of 200 draws\DIFdelbegin \DIFdel{, and the band of the nearest setting with a sibling grows from 0.25 to 2.9 floors}\DIFdelend . Selected on one geometry\DIFdelbegin \DIFdel{and evaluated on the other}\DIFdelend , the most decisive rule is \DIFdelbegin \DIFdel{the nearest setting with a sibling (2.0 and 3.6 floors ), the force trend or the median interpolating (4.2--5.8) and the force trend or the nearest setting extrapolating (8.9--11)}\DIFdelend \DIFaddbegin \DIFadd{a production-record rule in every relation within a family and decides at 2.0--11 floors on the other}\DIFaddend ; for a new family the rule selected on one transfer decides at 32--33 floors on the other\DIFdelbegin \DIFdel{, and the safest pre-specified rule , M3, is safe at 17--18}\DIFdelend \DIFaddbegin \DIFadd{. Selected for safety, the rule chosen on one geometry is safe on the other only at up to 18 floors interpolating and 11 extrapolating, and rule and band selected together keep 0.92--0.97 of the decided verdicts correct within a family, below 0.95 in four of six cases, and 0.73 for a new family (\ORt{15})}\DIFaddend . Paired intervals between point rules span about $\pm2$ floors with a sibling and when extrapolating, so smaller differences are not resolved by this design.

Requiring all seven characteristics of a design to meet raises the decisive distances of the production-record rules by at most 1.8 floors within a family \DIFdelbegin \DIFdel{(the nearest setting 5.1 floors with a sibling, the force trend 9.5 extrapolating) }\DIFdelend but those of M5 by up to \DIFdelbegin \DIFdel{16, and to 39 floors for M5 for a new family. Lower limitson the 5th percentile, calibration and truths }\DIFdelend \DIFaddbegin \DIFadd{16. Lower limits, truths and calibration }\DIFaddend from the first 20 or 50 parts, and \DIFdelbegin \DIFdel{the decided statistic $\bar m+4\sigma$ (an index of 1.33) }\DIFdelend \DIFaddbegin \DIFadd{further variants of data, constants and specifications }\DIFaddend leave the nearest setting within \DIFdelbegin \DIFdel{0.4 floor }\DIFdelend \DIFaddbegin \DIFadd{1.2 floors }\DIFaddend of its reference distance \DIFdelbegin \DIFdel{with a sibling and within 1.2 when extrapolating, }\DIFdelend and a production-record rule the most decisive within the family\DIFdelbegin \DIFdel{, as do variants of data, constants and specifications}\DIFdelend ; other floor protocols rescale the distances by 0.60 to 1.42, and floors recomputed without the held-out alternative's own series move the distances of the \DIFdelbegin \DIFdel{nearest setting and the force trend by at most 0.25 floors and those of the other }\DIFdelend main rules by at most 1.4 \DIFaddbegin \DIFadd{floors, that of M0 by 3.75}\DIFaddend .

%% ----------------------------------------------------------------------
\section{Discussion}\label{sec:discussion}

\subsection{What the evaluation establishes}

Two findings follow from the definitions and hold wherever the framework is applied: the floor puts every rule and characteristic on a common production-referenced scale, though not on a common scale of nonconforming fractions, and as a within-run quantity it bounds the variation of series production from below (RQ1); and an interval rule trades decisiveness for safety, which the two distances make explicit. The others are observations on one laboratory process, with every alternative produced once and a generic simulation dataset, and hypotheses for confirmation. The attainable distance grows with the novelty of the design: a production-record rule \DIFdelbegin \DIFdel{committed to in advance }\DIFdelend \DIFaddbegin \DIFadd{applied to every relation }\DIFaddend decided at 3.4 floors with a sibling\DIFdelbegin \DIFdel{(1.6 when the two produced siblings did not differ resolvably, 4.4 when they did)}\DIFdelend , 8.5 for a new setting \DIFdelbegin \DIFdel{(4.7 interpolating, 9.5 extrapolating) }\DIFdelend and 30 floors at best for a new family, and within a family the choice of rule mattered more than the relation, also among the pre-specified rules (RQ2). For this simulation dataset and for Gaussian processes and boosting fitted to six to nine alternatives, the simulation contributed little within a family, helping single characteristics when extrapolating in force, and mattered across families, for the characteristics it models directly; the learners added no decisiveness, and what made rules safe for new settings was \DIFdelbegin \DIFdel{an interval calibrated on the }\DIFdelend \DIFaddbegin \DIFadd{evidence borrowed from the other family on the same }\DIFaddend relation (RQ3). Qualification by relation also exposed a \DIFdelbegin \DIFdel{dangerous extrapolation}\DIFdelend \DIFaddbegin \DIFadd{danger}\DIFaddend : the scaled simulation, decisive with a sibling, accepted failing alternatives \DIFdelbegin \DIFdel{up to }\DIFdelend 165 floors from their limit when extrapolating\DIFdelbegin \DIFdel{.
}\DIFdelend \DIFaddbegin \DIFadd{, the multi-fidelity model up to 118.
}\DIFaddend 

%% guidance table (static)
\begin{table}[t]
\tabsize
\setlength{\tabcolsep}{4pt}
\caption{Guidance for design-stage manufacturability decisions, with its basis}\label{tab:guidelines}
\DIFdelbeginFL %DIFDELCMD < \begin{tabular}{@{}L{2.9cm}L{7.75cm}L{1.85cm}@{}}
%DIFDELCMD < \toprule
%DIFDELCMD < Situation & Guidance & Basis \\
%DIFDELCMD < \midrule
%DIFDELCMD < \multicolumn{3}{@{}l}{\emph{From the framework (prescriptive; usefulness not evaluated)}} \\
%DIFDELCMD < Before relying on a predictor & Measure the floor on consecutive production of the family with a stated protocol; report $\sigma_\mathrm{b}$, $\sigma_\mathrm{w}$, its lag dependence and the floors of single runs; treat it as a lower bound of series variation. & framework \\
%DIFDELCMD < Qualifying a rule & Evaluate it by grouped cross-validation for the evidence relation of the coming designs, per family, with siblings stratified by whether they differ from each other, and per characteristic. & framework \\
%DIFDELCMD < Choosing a rule & Compare it with the nearest produced setting and with its own version without the simulation, on both distances and by expected cost at the team's trial cost. & framework \\
%DIFDELCMD < Deciding a single design & Derive the guard band for the requirements the team expects, validate it out of sample, and apply it to the observable margin; a margin inside it means a trial. & framework \\
%DIFDELCMD < \multicolumn{3}{@{}l}{\emph{Observed in this study (one process, 18 alternatives produced once, 2 geometries; not general)}} \\
%DIFDELCMD < New variant with a produced sibling & the nearest produced setting decided beyond 3.4 floors (1.6 when the produced siblings did not differ resolvably from each other); no learned rule was more decisive & 18 alternatives \\
%DIFDELCMD < New process setting & production record 8.2--8.5 floors, interval rules 15--19; intervals matched to the relation made it safe & 6 force cells \\
%DIFDELCMD < New family & only rules with the (generic) simulation decided, at 30 floors and beyond; none qualified for single verdicts & two transfers \\
%DIFDELCMD < Release-level requirements & within about 2 floors of the 95th percentile: decided at best with a near-replicate sibling, otherwise by trial & 18 alternatives \\
%DIFDELCMD < \botrule
%DIFDELCMD < \end{tabular}
%DIFDELCMD < %%%
\DIFdelendFL \DIFaddbeginFL \begin{tabular}{@{}L{2.9cm}L{7.75cm}L{1.85cm}@{}}
\toprule
Situation & Guidance & Basis \\
\midrule
\multicolumn{3}{@{}l}{\emph{From the framework (prescriptive; usefulness not evaluated)}} \\
Before relying on a predictor & Measure the floor on consecutive production with a stated protocol; report $\sigma_\mathrm{b}$, $\sigma_\mathrm{w}$, lag dependence and single-run floors; treat it as a lower bound of series variation. & framework \\
Qualifying a rule & Evaluate it by grouped cross-validation for the relation of the coming designs, per family and characteristic, siblings stratified by whether they differ. & framework \\
Choosing a rule & Compare it with the nearest produced setting and with its own version without the simulation, on both distances and by expected cost at the team's trial cost. & framework \\
Deciding a single design & Take the guard band from the relations the family's production exercises or from another family, validate it out of sample for the expected requirements and widen it where trials are cheap; a margin inside it means a trial. & framework \\
\multicolumn{3}{@{}l}{\emph{Observed in this study (one process, 18 alternatives produced once, 2 geometries; not general)}} \\
New variant with a produced sibling & the nearest produced setting decided beyond 3.4 floors (1.6 when the produced siblings did not differ resolvably from each other); no learned rule was more decisive & 18 alternatives \\
New process setting & production record 8.2--8.5 floors, interval rules 15--19; borrowing the other family's errors on the relation made it safe & 6 force cells \\
New family & only rules with the (generic) simulation decided, at 30 floors and beyond; none qualified for single verdicts & two transfers \\
Release-level requirements & a limit at $P_\mathrm{pk}=1.33$ lies about 2 floors above the 95th percentile; confirming the index needed replicated runs, which neither a prediction nor a short trial replaced & 18 alternatives \\
\botrule
\end{tabular}
\DIFaddendFL \end{table}

\subsection{Implications for practice}

\DIFdelbegin \DIFdel{For a process planner or quality engineer the results say three things. }\DIFdelend Within a family the production record answers most design-stage questions \DIFaddbegin \DIFadd{of a process planner or quality engineer }\DIFaddend if it is indexed by how the coming design relates to it: a sibling at the same setting settles a requirement a few floors away, a new force within the produced range almost as well, and an extrapolated force or a new family needs a trial unless the requirement is far. Under the reference prior, of 100 verdicts on new variants with a sibling the nearest setting with its band sends 2 to trial and gets \DIFdelbegin \DIFdel{at least 93 }\DIFdelend \DIFaddbegin \DIFadd{92 to 95 }\DIFaddend of the other 98 right; of 100 on extrapolated settings it sends \DIFdelbegin \DIFdel{55 to trial. }\DIFdelend \DIFaddbegin \DIFadd{58 to trial, and under the release-like prior 41 and 92. }\DIFaddend And near release-level limits \DIFdelbegin \DIFdel{nothing replaced a trial without a nearly identical produced sibling}\DIFdelend \DIFaddbegin \DIFadd{neither a prediction nor a short trial confirmed capability, which needed replicated runs}\DIFaddend , so the value of design-stage prediction lies in the \DIFdelbegin \DIFdel{many }\DIFdelend decisions far from the limits\DIFdelbegin \DIFdel{rather than in the few close ones}\DIFdelend . For a developer of an AI-based manufacturability tool the lesson is methodological: a learned model has to be compared with a production-record baseline and with its own version without the simulation, scored by evidence relation in a unit that production defines, its intervals calibrated on cases of the same relation and shown to its user with a guard band. The nearest setting is also the most transparent rule, pointing to the produced alternatives \DIFdelbegin \DIFdel{whose record }\DIFdelend it averages, \DIFdelbegin \DIFdel{and so }\DIFdelend the natural default where a learned rule shows no advantage.

\subsection{Simulation and machine learning in this regime}

Within a design family the production record carried the decision even against a \DIFdelbegin \DIFdel{calibrated }\DIFdelend \DIFaddbegin \DIFadd{friction-fitted (censored) }\DIFaddend or rescaled simulation, but the simulation here is a generic dataset that does not contain the produced process \DIFdelbegin \DIFdel{: its force trend is two to five times too large, its friction calibration is censored at an edge of the grid with or without a discrepancy, and its oil-to-friction mapping is too sensitive}\DIFdelend \DIFaddbegin \DIFadd{(Section~\ref{sec:res-sim})}\DIFaddend . A process-calibrated simulation, with friction depending on pressure \citep{hol2012advanced}, a hardening model fit for springback \citep{wagoner2013advanced} and elastic tools \citep{pilthammar2018introduction}, might contribute more; the per-characteristic results show where it already helps, for the draw-ins and the arm angle when extrapolating and for the characteristics it models directly in a new family. The learners failed for a reason\DIFdelbegin \DIFdel{the framework exposes}\DIFdelend : with six to nine alternatives at two or three \DIFdelbegin \DIFdel{force levels }\DIFdelend \DIFaddbegin \DIFadd{forces }\DIFaddend a flexible model can identify at most a linear trend, and a well-specified one \DIFdelbegin \DIFdel{then }\DIFdelend becomes that trend. Nonlinear multi-fidelity models \citep{perdikaris2017nonlinear}, physics-informed learning \citep{karniadakis2021physics}, conformalised quantile regression on the parts \citep{romano2019conformalized} and multi-task models over more families \citep{bonilla2008multi} are the next candidates, to be scored in the same way.

\subsection{Reuse for other processes and AI-based tools: a proposal}\label{sec:reuse}

The framework needs several produced variants of a design family whose relations span those of the coming designs, a series of consecutive parts or replicated batches per variant, a design-stage predictor for every variant and characteristics measured identically on parts and predictions; here nine produced alternatives per family and 500-part series gave ten batches of 50 parts for the floor, one produced sibling sufficed for the nearest setting, and truths from the first 20 or 50 parts moved its distance with a sibling by less than half a floor. The extensions below are untested. A classifier with a fixed threshold does not take the requirement as input, so its operating characteristic is estimated by binning test cases by their true margin, which needs a continuous, production-referenced outcome for every test design that CAD-based manufacturability datasets usually lack. A language-model check is scored on a grid of requirements with repeated queries, reporting model\DIFdelbegin \DIFdel{and version, query }\DIFdelend \DIFaddbegin \DIFadd{, version, }\DIFaddend date, decoding settings and prompt paraphrases, treating refusals and hedged answers as trials \DIFdelbegin \DIFdel{, checking that its verdicts are monotone }\DIFdelend \DIFaddbegin \DIFadd{and testing its verdicts for monotonicity }\DIFaddend in the requirement and \DIFdelbegin \DIFdel{testing for contamination , since public datasets and the papers describing them may be in its training }\DIFdelend \DIFaddbegin \DIFadd{for contamination by the public }\DIFaddend data. A requirement on a rate, such as the share of split \DIFdelbegin \DIFdel{or wrinkled }\DIFdelend parts, is scored \DIFdelbegin \DIFdel{in the same way }\DIFdelend once each alternative's rate is estimated from its series, with a floor over batches of rates; for continuous design spaces the relation becomes a distance from the produced designs in descriptor space. Table~\ref{tab:checklist} is a proposed reporting checklist, modelled on model cards \citep{mitchell2019model} and TRIPOD+AI \citep{collins2024tripod}; the scoring is released as a small documented function \DIFdelbegin \DIFdel{, separate from the pipeline, }\DIFdelend that takes intervals, truths and floors and returns distances, guard bands and trial rates \citep{canseven2026archive}. An attempt on open powder-bed fusion data \citep{leirmo2021data} failed for \DIFdelbegin \DIFdel{lack of replicated batches within a build: a floor within a build would need at least }\DIFdelend \DIFaddbegin \DIFadd{want of }\DIFaddend ten replicated anchors of one design per build\DIFdelbegin \DIFdel{and several builds per setting}\DIFdelend .

\begin{table}[t]
\tabsize
\setlength{\tabcolsep}{3pt}
\caption{Proposed reporting checklist for evaluating a design-stage decision rule against production}\label{tab:checklist}
\DIFdelbeginFL %DIFDELCMD < \begin{tabular}{@{}L{3.3cm}L{4.6cm}L{4.4cm}@{}}
%DIFDELCMD < \toprule
%DIFDELCMD < Item & Why it matters & In this paper \\
%DIFDELCMD < \midrule
%DIFDELCMD < Floor and its protocol & every distance is in this unit & $B$, $\alpha$, lag, single runs (Table~\ref{tab:floors}); $\sigma_\mathrm{b}$, $\sigma_\mathrm{w}$ archived \\
%DIFDELCMD < Measurement of the characteristics & biases enter absolute comparisons & Table~\ref{tab:qcs}, Section~\ref{sec:qcs} \\
%DIFDELCMD < Relations, cases and replication units & alternatives, not cases, replicate & per relation, 18 alternatives \\
%DIFDELCMD < Both distances, false-accept side & decisiveness and safety differ & Table~\ref{tab:distances} \\
%DIFDELCMD < Baselines & production record, no simulation & NN, M1sn; every pair with and without \\
%DIFDELCMD < Intervals: source and level & coverage is a property of the construction & Table~\ref{tab:rules}, Section~\ref{sec:res-ml} \\
%DIFDELCMD < Per-characteristic results & pooled distances hide the worst & Table~\ref{tab:groups} \\
%DIFDELCMD < Guard band, prior, validation & the deployed part & per family, out of sample (Table~\ref{tab:guards}) \\
%DIFDELCMD < Costs and trial rates & deployment is a cost question & Fig.~\ref{fig:costs} \\
%DIFDELCMD < Development record & selection optimism & Section~\ref{sec:computation}, cross-fitting \\
%DIFDELCMD < Query protocol (language models) & verdicts vary between queries & proposal (Section~\ref{sec:reuse}) \\
%DIFDELCMD < \botrule
%DIFDELCMD < \end{tabular}
%DIFDELCMD < %%%
\DIFdelendFL \DIFaddbeginFL \begin{tabular}{@{}L{2.7cm}L{4.7cm}L{5.0cm}@{}}
\toprule
Item & Why it matters & In this paper \\
\midrule
Floor and protocol & every distance is in this unit & $B$, $\alpha$, lag, single runs (Tables~\ref{tab:floors}, S4) \\
Measurement & biases enter absolute comparisons & Table~\ref{tab:qcs}, Section~\ref{sec:qcs} \\
Relations, replication & alternatives, not cases, replicate & per relation, 18 alternatives \\
Both distances & decisiveness and safety differ & Table~\ref{tab:distances}, false-accept side \\
Baselines & production record, no simulation & NN, M1sn; rules with and without \\
Intervals & coverage follows the construction & Table~\ref{tab:rules}, Section~\ref{sec:res-ml} \\
Per characteristic & pooled distances hide the worst & Table~\ref{tab:groups} \\
Guard band, prior & the deployed part & decision-time, two priors (Table~\ref{tab:guards}) \\
Costs and trials & deployment is a cost question & Fig.~\ref{fig:costs} \\
Development record & selection optimism & Section~\ref{sec:computation}, cross-fitting \\
Query protocol & verdicts vary between queries & language models (Section~\ref{sec:reuse}) \\
\botrule
\end{tabular}
\DIFaddendFL \end{table}

\subsection{Deployment}

In use, a step-6 verdict would be shown with its margin in floors and in the unit of the characteristic, the relation on which it rests and the produced alternatives behind it, and a trial verdict with the requirement distance at which the rule would decide. A qualification has to be maintained: the realised errors of step-6 verdicts can be tracked against later production on a control chart of decision outcomes, \DIFaddbegin \DIFadd{which sees false accepts only, rejected designs never being made, }\DIFaddend and a rule re-qualified when its share of wrong decided verdicts exceeds 5\,\% beyond what the number of verdicts explains \DIFaddbegin \DIFadd{(a rise to 10\,\% takes about 150 decided verdicts to detect at 5\,\% significance and 80\,\% power)}\DIFaddend , when the floor of new runs leaves the range of the single-run floors, or when the die is reworked. The evidence has a cost, several 500-part series per family before any rule qualifies, \DIFdelbegin \DIFdel{which the results bound from below: }\DIFdelend \DIFaddbegin \DIFadd{though }\DIFaddend one produced sibling settles most verdicts with a sibling\DIFdelbegin \DIFdel{, and one produced alternative of a new family brings its decisive distance from 28 to 9 floors in its own floor}\DIFdelend . Whether teams decide better with the framework is an empirical question: a study with cases of known production outcome could test whether \DIFdelbegin \DIFdel{designers }\DIFdelend \DIFaddbegin \DIFadd{process planners and quality engineers }\DIFaddend shown the guard-banded verdict accept fewer wrong verdicts within the safe distance than \DIFdelbegin \DIFdel{designers }\DIFdelend \DIFaddbegin \DIFadd{those }\DIFaddend shown the verdict alone, at an equal rate of right acceptances \citep{lee2004trust}, \DIFaddbegin \DIFadd{how often they override a trial verdict, }\DIFaddend and whether the band narrows the spread of that share between teams \citep{zhang2021cautionary}.

\subsection{Threats to validity}

The evaluation covers one material, one press and two geometries, with relations within a family from a $3\times3$ factorial in which lubrication is mostly not resolved; the procedure is general, the numbers are not. The new family rests on two transfers between geometries that differ in three features at once. Each alternative was produced once, so the reference $q_p$ is that of one run\DIFdelbegin \DIFdel{and }\DIFdelend \DIFaddbegin \DIFadd{, without }\DIFaddend the variation between days, coils and tool states\DIFdelbegin \DIFdel{is missing}\DIFdelend ; the sensitivity analysis \DIFdelbegin \DIFdel{bounds }\DIFdelend \DIFaddbegin \DIFadd{explores }\DIFaddend its effect but cannot replace replicated runs. The floor does not separate scanner drift from process drift, and absolute values are uncalibrated. The simulation dataset was not calibrated to the process, so the findings on simulation concern this dataset. The requirements are one-sided limits on the 95th percentile of dimensional characteristics, below the conformance of series release; attribute outcomes such as splits and wrinkles are not covered; and the verdicts are marginal per characteristic, a design-level verdict being an extension. The guard band was estimated pooled over characteristics \DIFdelbegin \DIFdel{, }\DIFdelend \DIFaddbegin \DIFadd{and, for new settings, borrowed between two families; }\DIFaddend its prior is a reference rather than a team's, and the costs are assumed. Rules, strata, bands and priors were partly chosen after results \DIFdelbegin \DIFdel{had been seen and are }\DIFdelend \DIFaddbegin \DIFadd{were seen and }\DIFaddend evaluated on the same data; the cross-fitting \DIFdelbegin \DIFdel{bounds }\DIFdelend \DIFaddbegin \DIFadd{estimates }\DIFaddend the optimism but does not remove it, and a single analyst \DIFdelbegin \DIFdel{carried out }\DIFdelend \DIFaddbegin \DIFadd{did }\DIFaddend the analysis, with code and data open for replication.

%% ----------------------------------------------------------------------
\section{Conclusions}\label{sec:conclusion}

Design-stage manufacturability decisions can be evaluated against production on a scale production defines: the decisive and safe distances in floors tell how close to a requirement a rule decides and how close it can be relied on not to mislead, and a guard band validated for \DIFdelbegin \DIFdel{the requirements a team expects }\DIFdelend \DIFaddbegin \DIFadd{an assumed prior of requirements }\DIFaddend turns them into a rule for single decisions, priced by the costs of errors and trials. In an exploratory application to one laboratory deep-drawing process, with alternatives produced once and a generic simulation dataset, the attainable distance grew with the novelty of the design: the nearest produced setting, \DIFdelbegin \DIFdel{committed to in advance}\DIFdelend \DIFaddbegin \DIFadd{one rule for every relation}\DIFaddend , decided at 3.4 floors with a produced sibling and 8.5 for a new process setting, and no rule decided closer than 30 floors for a new family (the pre-specified M5 at 35). Release-level requirements lay within these distances, \DIFdelbegin \DIFdel{so they needed a trial unless a near-replicate sibling was produced}\DIFdelend \DIFaddbegin \DIFadd{and confirming a performance index of 1.33 needed replicated runs, which neither a prediction nor a short trial replaced}\DIFaddend .

For design practice, the production record carried the decision within a family, which argues for indexing such records by their relation to the coming design rather than for a more elaborate model; across families only rules using the simulation carried the change of geometry. For AI-based manufacturability support, the evaluation separated what the learners were credited with: Gaussian processes and boosting fitted to six to nine alternatives reproduced the production record, and \DIFdelbegin \DIFdel{intervals calibrated }\DIFdelend \DIFaddbegin \DIFadd{evidence borrowed from the other family }\DIFaddend on the same \DIFdelbegin \DIFdel{evidence }\DIFdelend relation, not learning, made it safe for new settings.

The evidence is one process and two geometries, and the comparison of rules is exploratory. Four extensions follow: series from several days, coils and tool states\DIFaddbegin \DIFadd{, scanned in a documented order with a reference artefact}\DIFaddend ; a process-calibrated simulation; families with continuous geometric descriptors; and a confirmatory test of the \DIFdelbegin \DIFdel{frozen procedure}\DIFdelend \DIFaddbegin \DIFadd{procedure, registered before the analysis, }\DIFaddend on other data. Within a design family, the baseline that a learned manufacturability check has to beat is a production record indexed by its relation to the coming design.


\backmatter

\DIFaddbegin \bmhead{Supplementary information}
\DIFadd{Online Resource~1 holds the tables behind statements of the paper (Tables~S1--S15), each with its result files.
}

\DIFaddend \bmhead{Acknowledgements}
The author thanks the authors of the RDDAC and DDACS datasets at the University of Stuttgart for publishing the data on which this work rests.

\section*{Declarations}

\bmhead{Funding}
No specific funding was received for this work.

\bmhead{Competing interests}
The author declares no competing interests.

\bmhead{Ethics approval and consent}
Not applicable.

\bmhead{Data availability}
All data are public. RDDAC: doi:10.18419/DARUS-5589 (version 2.0); DDACS: doi:10.18419/DARUS-4801 (version 3.0); both licensed CC BY 4.0. The extracted per-part and per-simulation feature tables and all result files are archived with the code \citep{canseven2026archive}.

\bmhead{Code availability}
The complete analysis pipeline, which reproduces every number in this paper from the public data, is archived on Zenodo \citep{canseven2026archive}; \texttt{results/findings.md} maps every number to the script and result file that produce it, \DIFaddbegin \texttt{\DIFadd{results/DATA\_DICTIONARY.md}} \DIFadd{translates its column and code names into the terms of the paper, }\DIFaddend and \texttt{scripts/opchar.py} is the scoring function for other datasets\DIFdelbegin \DIFdel{; M3n is named M4 in the code}\DIFdelend . A working copy is at \url{https://github.com/htcanseven/Claude}.

\bmhead{Author contributions}
H.T.C. conceived the method, designed and carried out the analysis, and wrote the paper.

\bmhead{Use of generative AI}
A large language model (Claude, Anthropic) was used under the author's direction to write and debug code, retrieve and check literature, run analyses and draft text. It was also used to simulate reviews of earlier versions of the manuscript in \DIFdelbegin \DIFdel{five }\DIFdelend \DIFaddbegin \DIFadd{six }\DIFaddend rounds, acting as referees and editor; most rules and analyses added after the first results, listed in \DIFdelbegin \DIFdel{Section~\ref{sec:computation}}\DIFdelend \DIFaddbegin \DIFadd{Online Resource~1}\DIFaddend , were added in response to these simulated reviews. None of these reviews is a review by a journal. The author designed the study, verified every analysis against the archived results, checked every reference against its published record, and takes full responsibility for the content; the language model is not an author. All figures come from the released code; none is AI-generated.


%% ----------------------------------------------------------------------
\FloatBarrier
{\small
\bibliography{refs}}

\end{document}
